\documentclass[10pt,reqno]{amsart}
\usepackage[T1]{fontenc}
\usepackage{lmodern}
\usepackage{microtype}
\usepackage{amsmath,amssymb,mathtools}
\usepackage{booktabs}
\usepackage[colorlinks=true,linkcolor=blue,citecolor=blue,urlcolor=blue]{hyperref}
\newtheorem{theorem}{Theorem}[section]
\newtheorem{lemma}[theorem]{Lemma}
\newtheorem{proposition}[theorem]{Proposition}
\newtheorem{corollary}[theorem]{Corollary}
\newtheorem{conjecture}[theorem]{Conjecture}
\newtheorem{introthm}{Theorem}
\renewcommand{\theintrothm}{\Alph{introthm}}
\theoremstyle{definition}
\newtheorem{definition}[theorem]{Definition}
\newtheorem{hypothesis}[theorem]{Hypothesis}
\newtheorem{remark}[theorem]{Remark}
\newtheorem{assessment}[theorem]{Assessment}
\newtheorem{problem}{Problem}
\numberwithin{equation}{section}
\newcommand{\Z}{\mathbb Z}
\newcommand{\N}{\mathbb N}
\newcommand{\Q}{\mathbb Q}
\newcommand{\R}{\mathbb R}
\newcommand{\C}{\mathbb C}
\newcommand{\HH}{\mathbb H}
\newcommand{\cL}{\mathcal L}
\newcommand{\cF}{\mathcal F}
\newcommand{\cS}{\mathcal S}
\newcommand{\Gam}{\Gamma}
\newcommand{\ind}{\mathbf 1}
\newcommand{\e}{\mathrm e}
\DeclareMathOperator{\SL}{SL}
\DeclareMathOperator{\PSL}{PSL}
\DeclareMathOperator{\vol}{vol}
\DeclareMathOperator{\disc}{disc}
\DeclareMathOperator{\lcm}{lcm}
\newcommand{\status}[1]{\textup{\textsc{#1}}}
\title[The Elsholtz--Tao Type I sum under Selberg's conjecture]{The Elsholtz--Tao Type~I sum under Selberg's eigenvalue conjecture}
\author{Anonymous}
\date{Draft, October 2026; internal checks only, not externally refereed}
\emergencystretch=3em
\begin{document}
\begin{abstract}
Assuming Selberg's eigenvalue conjecture for $\Gamma_0(4dq^2)$ with
even nebentypus modulo $2q$, uniformly for $d\ge1$ and odd squarefree
$q$ coprime to $d$, we prove $\sum_{p\le N}f_I(p)\ll N\log^2N$, removing the
$\log\log N$ from Elsholtz--Tao's upper bound for Type~I solutions of
$4/p=1/x+1/y+1/z$. The associated CM lifts of discriminant $-4d$ are
uniformly separated in $\HH$ ($\cosh\operatorname{dist}\ge3/2$).
Sobolev duality and the Deshouillers--Iwaniec/Drappeau spectral large
sieve give a per-$d$ count complementary to a per-$a$ Poisson--Weil count.
Their savings degenerate linearly at $d=a$, without costing $\log\log N$.
The proof is \emph{conditional} on the stated spectral hypothesis and
relative to cited inputs, including the inherited Elsholtz--Tao/internal-note
reduction. It has received two rounds of internal review only, not external
refereeing; that review did not re-check the inherited reduction.
Exceptional-spectrum estimates leave a positive-width strip untreated,
so this argument gives no unconditional improvement.
\end{abstract}
\maketitle
\setcounter{tocdepth}{1}
\tableofcontents

%%%%% part A
\section{Introduction}\label{sec:intro}
The Erd\H{o}s--Straus equation is
\begin{equation}\label{A:eq:es}
 \frac4n=\frac1x+\frac1y+\frac1z,\qquad x,y,z\in\N.
\end{equation}
Following Elsholtz--Tao \cite{ET}, a solution is of Type~I if $n\mid x$ and
$(n,yz)=1$, and of Type~II if $n\mid y,z$ and $(n,x)=1$.
The counts retain coordinate order; for odd primes,
$f(p)=3f_I(p)+3f_{II}(p)$ \cite[(1.4)]{ET}.
Their Theorem~1.1 states
\begin{equation}\label{A:eq:et}
 N\log^2N\ll\sum_{p\le N}f_I(p)\ll N\log^2N\log\log N,
 \qquad \sum_{p\le N}f_{II}(p)\asymp N\log^2N.
\end{equation}
On p.~5 of arXiv:1107.1010v6, after Theorem~1.1, they attribute the double
logarithm to Brun--Titchmarsh in very short progressions and conjecture its
elimination. In \S9, p.~36, a short modulus avoids the loss for Type~II.
Their parenthesis nevertheless says no similar trick seems available in the
``Type II case'' [sic]; reading this as Type~I is the interpretation of
\cite[\S1]{MN3}, not the literal quotation.

Put $\Delta=y^2(\partial_x^2+\partial_y^2)$, so that $-\Delta$ is nonnegative.
Here is the precise spectral hypothesis used in this note.
\begin{hypothesis}[SEL]\label{hyp:sel}
For every $d\ge1$, every odd squarefree $q\ge1$ with $(q,d)=1$, and every
even Dirichlet character $\chi$ modulo $2q$ (equivalently modulo $q$),
$-\Delta$ on $L^2(\Gamma_0(4dq^2)\backslash\HH,\chi)$ has no eigenvalue in
$(0,1/4)$.
\end{hypothesis}
The hypothesis covers the whole family; estimates are uniform in levels
and characters. Constants are removed before applying the spectral estimate.
\begin{introthm}\label{thm:main}
Assume \textup{(SEL)}. Then
\[
 \sum_{p\le N}f_I(p)\ll N\log^2N.
\]
\end{introthm}
This is Theorem~\ref{L:8.1} below. It concerns an average upper bound, not
existence of an Erd\H{o}s--Straus representation for every prime.

The new input is uniform separation of the level-$d$ CM points of
discriminant $-4d$. Sobolev duality then pairs them with a mean-zero
Poincar\'e series at a cost equal to the square root of their number times
a second-order Sobolev norm, with no level loss. A spectral large sieve
controls that norm under (SEL). The fixed-$d$ count complements a
fixed-$a$ Weil count precisely at $d=a$; Lemma~\ref{L:1.1} sums the saving
even though it degenerates linearly at that boundary. Only absolute
errors averaged over $d$, not uniform per-packet relative errors, are used.

Unconditionally, the available estimates leave a positive-width strip
next to $d=a$; a direct exponent comparison gives a scale about $7/32$
in $2\alpha-1+\gamma$, not a sharp boundary. This argument gives
\emph{no unconditional improvement} of \eqref{A:eq:et}; see \S\ref{sec:uncond}.

\paragraph{Status and dependencies.}
Theorem~\ref{thm:main} is \status{conditional} on (SEL), with a proof relative
to cited inputs. These are the ET reduction (Proposition~2.2, Lemma~2.8,
and (8.1)--(8.2)); the internal note \cite[Proposition~2.3 and
proof of Theorem~3.8, step~(1)]{MN3} (MN3's Proposition~2.3 is
Proposition~\ref{M:2.3} here; R114 repair); Deshouillers--Iwaniec \cite[Theorem~2]{DI} and
Drappeau \cite[Proposition~4.7, arXiv numbering]{Dr}; the spectral theorem;
Weil's bound; Brun--Titchmarsh; Selberg's upper-bound sieve; and
P\'olya--Vinogradov. DI Theorem~5 enters the discussion of exceptional
spectrum, not the conditional proof. The new argument received internal
adversarial review (R111, rounds 1--2), not external refereeing. That review did
not re-check the inherited ET/MN3 reduction; MN3 received a separate internal
review. (R114 repair) A later internal referee pass (R114) re-derived the
parts of that reduction used here, namely $f_I(p)\le2\sum_cw_c(p)$ from
\cite[Proposition~2.2, Lemma~2.8]{ET} and the large-$c$ step from
\cite[(8.1)--(8.2)]{ET}; this is still internal checking only.
The book locators in \cite{Iw} are inherited from the source note
and have not been independently verified; the same applies to the
locators in \cite{HR} and to the precise statements of \cite{MV} and
\cite{KS}, which were not available to the internal reviewers (R114 repair). Computational checks are not proofs.

\paragraph{Literature and scope.}
A limited literature search found no later removal of the Type~I double
logarithm; it was not an exhaustive priority search. Jia \cite{Jia} is cited
in ET's discussion following Theorem~1.1, but for the Type~II prime sum; the
final arXiv version still calls the Type~I double logarithm a technical
limitation and conjectures its removal. Duke \cite{Duke} and Liu--Masri--Young \cite{LMY} address related CM
distributions, but do not supply the joint level--discriminant estimate
needed here, with level comparable to the discriminant.

\section{The Type I sum and the bad region}\label{sec:reduction}
Write $\cL=\log N$ and work first with $N/2<n\le N$. The nonnegative weight
\begin{equation}\label{A:eq:weight}
 w_c(n)=\#\{(a,d,f)\in\N^3:f\mid4a^2d+1,\ n=4acd-f,
                         \ n/4<acd\le3n/4\}
\end{equation}
satisfies $f_I(p)\le2\sum_c w_c(p)$, by \cite[Proposition~2.2,
Lemma~2.8]{ET}; the factor two restores the order of $y,z$.
ET's map is $(x,y,z)=(abdn,acd,bcd)$, with $(abcd,n)=1$ and
$\gcd(a,b,c)=1$. With $y\le z$, it has
\begin{equation}\label{A:eq:identities}
 4abd=ne+1,\quad ce=a+b,\quad4acd=n+f,\quad
 ef=4a^2d+1,\quad bf=na+c,
\end{equation}
and $a\le b$, $an\le bf\le5an/3$.
Conversely, for every tuple in \eqref{A:eq:weight}, put
$e=(4a^2d+1)/f$ and $b=ce-a$. Direct substitution gives all the identities
\eqref{A:eq:identities}, but only $b\ge a/2$, not necessarily $b\ge a$:
indeed $0<f\le2n$ and $bf=na+c$. We retain this larger set throughout.
Writing $a=N^\alpha$, $b=N^\beta$, $c=N^\gamma$, we have
\begin{equation}\label{A:eq:exponents}
 d\asymp N^{1-\alpha-\gamma},\quad e\asymp N^{\beta-\gamma},\quad
 f\asymp N^{1+\alpha-\beta},\quad
 \beta\ge\alpha-O(\cL^{-1}),\quad\alpha+\gamma\le1+O(\cL^{-1}).
\end{equation}
All exponent regions below are understood up to these bounded dyadic errors.

\begin{lemma}[SL$_2$ form; {\cite[Lemma~3.1]{MN3}}]\label{M:3.1}
The identity $ef-4a^2d=1$ is equivalent to
\[
 M=\begin{pmatrix}e&2a\\2ad&f\end{pmatrix}\in\SL_2(\Z),
 \qquad n=2cM_{21}-M_{22}.
\]
If $T=\operatorname{diag}(1,d)$, then
$\sigma_d(M)=TM^{\mathsf t}T^{-1}$ is an anti-involution and
$S_d=\{M\in\SL_2(\Z):M_{21}=dM_{12}\}$ is its fixed set.
It is stable under $M\mapsto\gamma M\sigma_d(\gamma)$ for
$\gamma\in\Gamma_0(d)$. The positive Type~I data also give the positive
definite form $[f,4ad,de]$ of discriminant $-4d$.
\end{lemma}
\begin{proof}
The determinant and linear-form assertions follow by multiplication.
For $\gamma=\left(\begin{smallmatrix}p&t\\r&s\end{smallmatrix}\right)$,
$\sigma_d(\gamma)=\left(\begin{smallmatrix}p&r/d\\dt&s\end{smallmatrix}\right)$
is integral exactly when $d\mid r$; reversing products proves stability.
The form has discriminant $16a^2d^2-4def=-4d$ and positive leading
coefficient. Conversely its coefficients recover $f,a,e$ when the middle
coefficient is divisible by $4d$. Positivity and this parity restriction
are extra conditions, not invariants of all of $S_d$.
\end{proof}

\begin{remark}[Single-progression parametrisations]\label{M:3.2}
The generic elimination in \cite[Lemma~3.2]{MN3} produces the moduli
\[
 4ad,\quad4bd,\quad4ab,\quad4acf,\quad4cdf,
 \qquad4bcf',\quad4cdf',\qquad f'=4bcd-n.
\]
Its claim that these exhaust all three-coordinate single-progression
parametrisations has status \status{evidence}: degenerate eliminants were
not checked systematically. No completeness claim is used here. The
following explicit fibres suffice. Fixing $(a,c,f)$ forces
$d\equiv-(4a^2)^{-1}\pmod f$ and puts $n$ in one class modulo $4acf$.
Fixing $(c,d,f)$ puts $a$ in
$\rho_d(f)=\#\{x\bmod f:4dx^2+1\equiv0\pmod f\}$ classes, hence $n$ in
that many classes modulo $4cdf$. Fixing $(a,b,c)$ fixes $e=(a+b)/c$;
then $(e,4ab)=1$, $d\equiv(4ab)^{-1}\pmod e$, and
$n=(4abd-1)/e$ lies in one class modulo $4ab$.
Every map from the free variable to $n$ is injective. For fixed $(a,b)$
there are at most $\tau(a+b)$ choices of $c$. The moduli $4ad,4bd$ arise by
varying $c$ with the appropriate divisor fixed; swapping $a,b$ gives the twins.
\end{remark}

\begin{proposition}[Bad-region geometry; {\cite[Proposition~3.3]{MN3}}]\label{M:3.3}
In the ideal exponent domain $0\le\alpha\le\beta$, $\alpha+\gamma\le1$,
all seven displayed moduli have exponent at least $1-\eta$ exactly on
\begin{equation}\label{A:eq:bad}
 \mathcal R_{\rm bad}(\eta)=
 \{\gamma\le\eta,\ \alpha+\beta\ge1-\eta,\
              \beta\le2\alpha+\gamma+\eta,\ \beta\le1+\eta\}.
\end{equation}
Its limiting $\gamma=\eta=0$ slice has area $1/6$. There
$e,f$ are at least $\max(a,d)$ in exponent; for positive $\eta$ this holds
only up to $N^{-O(\eta)}$.
\end{proposition}
\begin{proof}
The seven exponents, in the order of the preceding list, are
\[
 1-\gamma,\quad1-\gamma+\beta-\alpha,\quad\alpha+\beta,\quad
 1+2\alpha+\gamma-\beta,\quad2-\beta,\quad
 1+2\beta+\gamma-\alpha,\quad2-2\alpha+\beta.
\]
The second and last two inequalities are redundant, giving
\eqref{A:eq:bad}. At $\gamma=\eta=0$ the slice is
$\alpha+\beta\ge1$, $\alpha\le\beta\le\min(2\alpha,1)$, of area
$\int_{1/3}^{1/2}(3\alpha-1)\,d\alpha+
\int_{1/2}^1(1-\alpha)\,d\alpha=1/6$.
Both $\beta$ and $1+\alpha-\beta$ are at least
$\max(\alpha,1-\alpha)$ there. Perturbing the inequalities proves the last
assertion. For our enlarged weights $b\ge a/2$ introduces only
$O(\cL^{-1})$ errors.
\end{proof}
A positive area alone does not prove positive logarithmic mass or exclude
other parametrisations. We use the region only to organise an upper bound.

\subsection{Arithmetic estimates outside the bad region}
We record the imported divisor estimate explicitly. Here $\omega(k)$ is
the number of distinct prime divisors.
\begin{proposition}[Coprimality gain; {\cite[Proposition~2.3]{MN3}}]\label{M:2.3}
Fix $l\ge1$ and let $k\ge1$, $U,V\ge2$.
\textup{(a)} If $kV^2\le U^l$ and $U\ge\omega(k)+2$, then
\[
 \sum_{u\le U,v\le V}\tau(kuv^2+1)
       \ll_l\frac{\varphi(k)}kUV\log U.
\]
\textup{(b)} If $kU\le V^l$ and $V\ge\omega(2k)+2$, then the same sum is
\[
 \ll_l\frac{\varphi(k)}kUV\log V
 \left[1+\min\left\{\log(1+k),
              \frac{k^{1/6}\log(2kU)}{\sqrt U}\right\}\right].
\]
\end{proposition}
This is a proved input relative to \cite[Theorem~7.1, Corollary~7.4]{ET}
and P\'olya--Vinogradov. Its proof in \cite[\S2]{MN3} keeps the indicator
that divisors are coprime to $k$, rather than dropping it in the quadratic
root bound. We do not re-prove that imported proposition.

Put $g(x)=x/\varphi(x)$ and write $x\asymp X$ for a dyadic interval
$X<x\le2X$, allowing a separate initial interval containing $1$.
\begin{lemma}[Weighted divisor sums]\label{L:8.3}
Uniformly in $A,B,D,F,X\ge1$,
\begin{align}
 \textup{(a)}\quad\sum_{x\asymp X}g(x)^2&\ll X,\qquad
             g(x)=\sum_{l\mid x}\frac{\mu^2(l)}{\varphi(l)},\label{A:eq:gmean}\\
 \textup{(b)}\quad\sum_{\substack{a\asymp A,b\asymp B\\a\le2b}}
          \tau(a+b)g(a)g(b)&\ll AB\log(2B),\label{A:eq:abmass}\\
 \textup{(c)}\quad\sum_{d\asymp D}g(d)\sum_{f\asymp F}\frac{\rho_d(f)}{\varphi(f)}&\ll D.
                                                        \label{A:eq:rootmass}
\end{align}
\end{lemma}
\begin{proof}
The identity in \eqref{A:eq:gmean} follows on prime powers. Also
$g^2=1*h$ with $h\ge0$, supported on squarefree integers, and
$h(p)=(p/(p-1))^2-1=O(1/p)$. Thus $\sum_lh(l)/l<\infty$, which proves
the mean square and, by Cauchy--Schwarz, $\sum_{x\asymp X}g(x)\ll X$.
The same divisor expansion gives $\sum_{x\le X}1/\varphi(x)\ll\log(2X)$.

For \eqref{A:eq:abmass}, fix $a$ and retain $a\le2b$, so $a+b\le6B$.
The elementary bound $\tau(n)\le2\sum_{\delta\mid n,\,\delta\le\sqrt n}1$
gives, for squarefree $l$,
\[
 \sum_{\substack{b\asymp B,\ l\mid b\\a\le2b}}\tau(a+b)
 \le2\sum_{\delta\le\sqrt{6B}}
       \left(\frac{2B(l,\delta)}{l\delta}+1\right)
 \ll\frac B l\tau(l)\log(2B)+\sqrt B.
\]
Multiply by $\mu^2(l)/\varphi(l)$ and sum $l\le2B$.
The series $\sum_l\mu^2(l)\tau(l)/(l\varphi(l))$ converges and the
remaining sum is $O(\sqrt B\log(2B))$. Summing with weight $g(a)$ proves
\eqref{A:eq:abmass}.

For \eqref{A:eq:rootmass}, set $\chi_d=(\frac{-4d}{\cdot})$ and
$S_d(Y)=\sum_{r\le Y}\chi_d(r)/r$, with $S_d(Y)=0$ if $Y<1$.
This is a nonprincipal real Dirichlet character modulo $4d$.
The prime-power root counts give
\[
 \rho_d(f)\le(1*\chi_d)(f),\qquad
 \rho_d(kf)\le2^{\omega(k)}\rho_d(f)\quad(k\text{ squarefree}).
\]
Indeed even moduli have no roots; for odd $p\nmid d$ the count at every
positive power is $1+\chi_d(p)$, and for $p\mid d$ it is zero.
For $Y\ge1$,
\begin{equation}\label{A:eq:convolution}
 \sum_{Y<f\le2Y}(1*\chi_d)(f)=Y S_d(2Y)+O(Y).
\end{equation}
The floor-function error is $O(\sum_{r\le2Y}1)=O(Y)$.
We claim, uniformly in $Y\ge1$,
\begin{equation}\label{A:eq:PVmean}
 \sum_{d\asymp D}|S_d(Y)|^2\ll D.
\end{equation}
For $Y\le Y_0=D/\log^2(2D)$ expand the square. Terms with $rr'$ even
vanish. If $rr'$ is a square the inner $d$-sum is at most $O(D)$, and
$\sum_{rr'=\square}(rr')^{-1}<\infty$ (write $r=su^2,r'=sv^2$ with $s$
squarefree). Otherwise P\'olya--Vinogradov for
$d\mapsto(\frac d{rr'})$ bounds that sum by
$O(\sqrt{rr'}\log(2rr'))$. The total is
$O(D+Y\log(2Y))=O(D)$. For $Y>Y_0$, partial summation and
P\'olya--Vinogradov modulo $4d$ give
\[
 S_d(Y)-S_d(Y_0)\ll\frac{\sqrt D\log(2D)}{Y_0}=O(1).
\]
The last assertion is uniform for large $D$; bounded $D$ are handled
directly by periodicity and partial summation. This proves
\eqref{A:eq:PVmean}; see \cite[Chapter~12]{IK} for the character-sum input.
Finally the divisor expansion of $g(f)$, the root inequalities and
\eqref{A:eq:convolution} imply
\[
 \sum_{f\asymp F}\frac{\rho_d(f)}{\varphi(f)}
 \ll\sum_{k\ge1}\frac{\mu^2(k)2^{\omega(k)}}{k\varphi(k)}
                 (|S_d(2F/k)|+1).
\]
When $F/k<1$ the inner interval contains at most $1$, so the same bound
holds. Apply Cauchy--Schwarz using \eqref{A:eq:gmean} and
\eqref{A:eq:PVmean}, and sum the convergent $k$-series.
\end{proof}

\begin{lemma}[Brun--Titchmarsh outside the bad region]\label{L:8.2}
Fix $0<\eta_1<1/2$. For $P\in\{ab,acf,cdf\}$, the number of tuples
in \eqref{A:eq:weight}, with arbitrary $c$, $n$ prime in $(N/2,N]$, and
$m_P\le N^{1-\eta_1}$, is $O(\eta_1^{-1}N\cL^2)$, where
$m_{ab}=4ab$, $m_{acf}=4acf$, $m_{cdf}=4cdf$.
\end{lemma}
\begin{proof}
For $N\ge N_0(\eta_1)$, Brun--Titchmarsh \cite{MV} on an interval of
length $N/2$ gives $O(N/(\eta_1\cL\varphi(m)))$ primes per primitive
class. A nonprimitive class contains at most one prime. By the injective
fibres above, each fixed triple contributes at most
\[
 \rho_T\left(\frac{2N}{\eta_1\cL\varphi(m_T)}+1\right),
 \qquad \rho_T=1,1,\rho_d(f),
\]
respectively for $ab,acf,cdf$, with an inessential enlargement of the
absolute constant. Use $\varphi(4x)\ge2\varphi(x)$ and
$\varphi(xy)\ge\varphi(x)\varphi(y)$.
For $acf$ the reciprocal-totient sum is $O(\cL^3)$.
For $ab$ it is bounded by
$\sum_{ab\le N,a\le2b}\tau(a+b)/(\varphi(a)\varphi(b))=O(\cL^3)$,
by \eqref{A:eq:abmass} on $O(\cL^2)$ boxes.
For $cdf$, \eqref{A:eq:rootmass} gives
$\sum_{df\le N}\rho_d(f)/(\varphi(d)\varphi(f))=O(\cL^2)$;
the $c$-sum supplies one more logarithm.
The $+1$ errors total $O_\epsilon(N^{1-\eta_1+\epsilon})$, by the divisor
bound and $\rho_d(f)\le2^{\omega(f)}$. Choose $\epsilon<\eta_1/2$.
\end{proof}

\begin{lemma}[Weighted Type I masses]\label{L:8.4}
Let $\cL^2\le A\le N$ and $1\le D\le N$.
\textup{(a1)} If $D\le A$, then
$\sum_{a\asymp A,d\asymp D}\tau(4a^2d+1)g(d)\ll AD\cL$;
if $D\ge A$, the same holds with $g(a)$ in place of $g(d)$.
\textup{(a2)} If also $A,D\ge N^{1/4}$, then
\begin{equation}\label{A:eq:doublemass}
 \sum_{a\asymp A,d\asymp D}\tau(4a^2d+1)g(a)g(d)\ll AD\cL.
\end{equation}
\textup{(b)} For every $A,D\ge1$ and $F\le4A$,
\begin{equation}\label{A:eq:smallmass}
 \sum_{d\asymp D}g(d)\#\{(a,f):a\asymp A,f\asymp F,
                                    f\mid4a^2d+1\}\ll AD.
\end{equation}
The last assertion also holds for $e$ in place of $f$.
\end{lemma}
\begin{proof}
Expand $g(d)=\sum_{l\mid d}\mu^2(l)/\varphi(l)$ when $D\le A$ and put
$d=ld'$. Proposition~\ref{M:2.3}, second assertion, with coefficient $4l$,
linear range $2D/l$, quadratic range $2A$, and fixed exponent $40$, gives
\[
 \sum_{a\le2A,d'\le2D/l}\tau(4la^2d'+1)
       \ll\frac{AD}{l}\cL(1+\log(1+4l)).
\]
The required size inequalities are $8D\le(2A)^{40}$ and
$2A\ge\omega(8l)+2$. If a range is smaller than $2$, enlarge it to $2$;
then $8l\le(2A)^{40}$ suffices and the loss is bounded since $l\le2D$.
The resulting $l$-series with denominator $l\varphi(l)$ converges.
For $D\ge A$, expand $g(a)$, put $a=ka'$, and use the first assertion of
Proposition~\ref{M:2.3} with coefficient $4k^2$, ranges $2D,2A/k$.
Here $16A^2\le(2D)^{40}$ and $2D\ge\omega(4k^2)+2$; after enlarging a
small range use $16k^2\le(2D)^{40}$. The bound is
$O((AD/k)\cL)$, summable against $1/\varphi(k)$.

For \eqref{A:eq:doublemass}, expand both factors and write $a=ka'$,
$d=ld'$, with coefficient $K=4k^2l$. The part $k>A^{1/2}$ or
$l>D^{1/2}$ is, for a sufficiently small fixed $\epsilon>0$,
\[
 \ll N^\epsilon AD(A^{-1/2}+D^{-1/2})\ll AD.
\]
Here the divisor bound is uniform since $4a^2d+1\ll N^3$, and
$\sum_{k>T}1/(k\varphi(k))\ll T^{-1}$ follows from
\eqref{A:eq:gmean}. In the remaining part, if $D\ge A$ use the first
assertion of Proposition~\ref{M:2.3}: its coefficient condition is
$16lA^2\le(2D/l)^{40}$. If $D<A$ use the second assertion, whose condition
is $8k^2D\le(2A/k)^{40}$. These hold for all sufficiently large $N$,
since the relevant large variable is at least $N^{1/8}$ and the left
side is $O(N^3)$. The conditions involving $\omega(K)$ hold as well.
Discarding $\varphi(K)/K\le1$, both bounds are summable using
\[
 \sum_{k,l\ge1}\frac{\mu^2(k)\mu^2(l)(1+\log(1+4k^2l))}
                         {kl\varphi(k)\varphi(l)}<\infty.
\]
Finally a fixed $(d,f)$ admits at most $\rho_d(f)(A/f+1)$ values of $a$.
For $f\le2F\le8A$ this is at most $9A\rho_d(f)/f$.
Equation~\eqref{A:eq:rootmass} proves \eqref{A:eq:smallmass}.
\end{proof}

\subsection{What remains}
Only the following unconditional step of the earlier LD reduction is needed.
\begin{proposition}[Large-$c$ reduction; {\cite[proof of Theorem~3.8, step~(1)]{MN3}}]\label{M:3.8}
For fixed small $\eta>0$, the contribution of $c>N^\eta$ to
$\sum_{N/2<p\le N}\sum_c w_c(p)$ is $O(\eta^{-1}N\cL^2)$.
\end{proposition}
\begin{proof}
Here $ad\ll N^{1-\eta}$. Brun--Titchmarsh in
$n\equiv-f\pmod{4ad}$ and \cite[(8.1)--(8.2)]{ET} on product blocks give
\[
 \frac{N}{\eta\cL}\sum_{ad\ll N}\frac{\tau(4a^2d+1)}{\varphi(ad)}
       \ll\eta^{-1}N\cL^2.
\]
The classes are primitive since $f\mid4a^2d+1$ implies $(f,4ad)=1$,
also for the enlarged weights.
\end{proof}

It remains to count prime tuples with $c\le N^\eta$ and
\begin{equation}\label{A:eq:remaining}
 \min(4ab,4acf,4cdf)>N^{1-\eta_1}\quad(\eta_1\ge\eta).
\end{equation}
The moduli $4ad,4bd$ are already $\gg N^{1-\eta}$ because
$4acd\asymp N$ and $b\ge a/2$; the twin exponents are automatically large.
Thus \eqref{A:eq:remaining} lies in $\mathcal R_{\rm bad}(\eta_1)$ up to
$O(\cL^{-1})$. Smooth upper-bound cells may be supported in
$\mathcal R_{\rm bad}(2\eta_1)$, where
$e,f\ge N^{1/2-3\eta_1}$ up to absolute constants. This is the only
remaining region; in particular no assertion of completeness from
Remark~\ref{M:3.2} is needed.

\section{\texorpdfstring{SL$_2$}{SL2}/Heegner reformulation and separation}\label{sec:heegner}
For $d\ge1$ define
\[
 \cF_d=\{[A,B,C]\in\Z^3:B^2-4AC=-4d,\ A>0,
                         \ 2d\mid B,\ d\mid C\}.
\]
Its root map is $z_Q=(-B+i\sqrt{4d})/(2A)$.
Every form here is positive definite. Type~I forms form the subset
\[
 \cF_d^I=\{[f,4ad,de]:ef-4a^2d=1,\ e,f>0,\ a\in\Z\}.
\]
The unrestricted integer $a$ makes the invariant set convenient; a positive
box will restore $a>0$. These are CM, or Heegner, forms with discriminant
$-4d$ and level $d$, not forms satisfying the coprime Heegner hypothesis:
$(-4d,d)=d>1$ for $d>1$ (R114 repair). Type~I forms are primitive, since $(f,4ad)=1$.
The larger set $\cF_d$ need not be primitive (e.g.\ $[2,6,6]\in\cF_3$);
Lemma~\ref{L:6.3} counts only the Type~I subset.

\begin{lemma}[Distance formula]\label{L:2.1}
For positive definite forms $Q,Q'$ of equal discriminant $D<0$,
\[
 \cosh\operatorname{dist}_{\HH}(z_Q,z_{Q'})
             =1+\frac{\disc(Q-Q')}{2|D|}.
\]
\end{lemma}
\begin{proof}
Write $Q=[A,B,C]$, $Q'=[A',B',C']$. The formula
$\cosh\operatorname{dist}(z,w)-1=|z-w|^2/(2\Im z\Im w)$ gives
\[
 \cosh\operatorname{dist}(z_Q,z_{Q'})-1
 =\frac{(BA'-B'A)^2+|D|(A-A')^2}{2|D|AA'}.
\]
Substituting $4AC=B^2+|D|$ and $4A'C'=B'^2+|D|$ into the discriminant
of the difference gives
\[
 AA'\disc(Q-Q')=(BA'-B'A)^2+|D|(A-A')^2,
\]
which proves the assertion.
\end{proof}
\begin{lemma}[Uniform separation]\label{L:2.2}
Distinct $Q,Q'\in\cF_d$ satisfy
$\cosh\operatorname{dist}_{\HH}(z_Q,z_{Q'})\ge3/2$, independently of $d$.
\end{lemma}
\begin{proof}
We have $4d^2\mid(B-B')^2$ and $4d\mid4(A-A')(C-C')$, hence
$4d\mid\disc(Q-Q')$. By Lemma~\ref{L:2.1},
$\cosh\operatorname{dist}-1\in\frac12\Z_{\ge0}$.
If it is zero, the roots coincide. A positive definite form of given
negative discriminant is determined by its root: in
$Q=A(X-zY)(X-\bar zY)$ the discriminant fixes $A>0$.
Thus equality zero implies $Q=Q'$, which is excluded.
\end{proof}
Separation persists on every subset, including parity and sieve subsets.
It is a statement about all lifts in $\HH$, not just a chosen list of
representatives on a quotient. This distinction permits the uniform
Sobolev argument of \S\ref{sec:sobolev}.

\subsection{Coordinates, parity and the sieve group}
Set
\begin{equation}\label{A:eq:wcoord}
 w_Q=z_Q/d=\frac{-2a+i/\sqrt d}{f}\quad(Q\in\cF_d^I).
\end{equation}
Dilation is a hyperbolic isometry. The set $\cF_d$ is preserved by
$Q\mapsto Q\circ\gamma$ for $\gamma\in\Gamma^0(d)$, where
$\Gamma^0(d)$ means $d\mid\gamma_{12}$. Indeed
\begin{align*}
 C'&=At^2+Bts+Cs^2\equiv0\pmod d,\\
 B'&=2Apt+B(ps+tr)+2Crs\equiv0\pmod{2d}.
\end{align*}
for $\gamma=\left(\begin{smallmatrix}p&t\\r&s\end{smallmatrix}\right)$.
Roots transform by $\gamma^{-1}$. Conjugating by $z=dw$ identifies this
with $\Gamma_0(d)$ on the $w$-side.

The smaller set $\cF_d^I$ requires $4d\mid B$. On the $z$-side it is
preserved by $\Gamma^0(2d)\cap\Gamma_0(2)$, since each term in the
formula for $B'$ is then divisible by $4d$. On the $w$-side this becomes
$\Gamma_0(2d)\cap\Gamma^0(2)$.
In particular, for squarefree odd $(q,d)=1$, take
\begin{equation}\label{A:eq:sievegroup}
 \Gamma'_q=\Gamma_0(2d)\cap\Gamma(2q),\qquad u=w/(2q),\qquad M=4dq^2.
\end{equation}
This preserves both the parity and the condition
$q\mid n(Q)=cB-A$. To see the latter, the corresponding $z$-matrix is
congruent to $I$ modulo $q$, using $(q,d)=1$; hence it preserves the
coefficients of $Q$ modulo $q$. It is important not to replace this group
by $\Gamma_0(d)\cap\Gamma(2)$: for even $d$ that group need not preserve
$4d\mid B$ on the $w$-side.

In $u$-coordinates the effective group (adjoining $-I$, which acts trivially)
is
\[
 \Gamma_{M,q}=
 \left\{\pm\begin{pmatrix}p&t\\r&s\end{pmatrix}\in\Gamma_0(M):
                                     p\equiv s\equiv1\pmod{2q}\right\}.
\]
Its cusp at infinity has width one, and it contains $\Gamma_1(M)$.
For $q>1$, its index in $\Gamma_0(M)$ is $\varphi(q)/2$; for $q=1$
the group is simply $\Gamma_0(4d)$ and the index is $1$.
The quotient for $q>1$ is $(\Z/q\Z)^\times/\{\pm1\}$.
Consequently the spectral decomposition on $\Gamma_{M,q}$ splits into
the even-character spaces on $\Gamma_0(M)$ (with just the trivial space
when $q=1$). This explains precisely the family in (SEL).

\subsection{Boxes and Poincar\'e series}
For a smooth compactly supported weight in $a\asymp A,f\asymp F$, its
expression on $\HH$ is
\[
 \psi(w)=W_1\!\left(\frac{-\Re w}{2A\Im w\sqrt d}\right)
          W_2\!\left(\frac1{F\sqrt d\Im w}\right).
\]
Its height is $\asymp(F\sqrt d)^{-1}$, width $\asymp A/F$, and
hyperbolic area $\asymp A\sqrt d$. For spectral estimates use a partition in
$(\log f,\log(a/f))$. In the $u$-coordinate the cells then have the
product shape
\begin{equation}\label{A:eq:box}
 \psi(u)=\phi(x/\lambda)W(y/Y),\qquad
 \lambda\asymp\frac A{qF},\qquad Y\asymp\frac1{qF\sqrt d},
\end{equation}
where $\operatorname{supp}\phi\subset[-2,-1]$ and
$\operatorname{supp}W\subset[1,2]$.
Fixed smooth enlargements cover the original boxes by positivity. In
particular $Y/\lambda\asymp(A\sqrt d)^{-1}$. In the bad region this
is a negative power of $N$, so the condition $Y\le\lambda L_*^{-3}$
of Proposition~\ref{L:5.1} holds, since its local logarithm $L_*\ll\cL$.

Here is the counting identity with the stabilisers made explicit.
Let $\widetilde\Lambda$ be the set of points arising from an invariant
subset of $\cF_d$ in the chosen coordinate, let $\Gamma$ be its group,
and let its cusp width at infinity be $h$. Set
\[
 \psi_{\rm per}(z)=\sum_{m\in\Z}\psi(z+mh),\qquad
 P_\psi(z)=\sum_{\gamma\in\Gamma_\infty\backslash\Gamma}
                                     \psi_{\rm per}(\gamma z).
\]
Write $\Lambda=\Gamma\backslash\widetilde\Lambda$ and
$e_z=|\Gamma_z/(\Gamma\cap\{\pm I\})|$. This set is finite: there are
finitely many integral-form classes of fixed discriminant, and $\Gamma$
has finite index in the appropriate conjugate of $\SL_2(\Z)$.
With $d\mu=dx\,dy/y^2$ and $\#\Lambda=\sum_{z\in\Lambda}e_z^{-1}$,
unfolding gives
\begin{equation}\label{A:eq:countidentity}
 N_d(\psi):=\sum_{w\in\widetilde\Lambda}\psi(w)
 =\sum_{w\in\Gamma_\infty\backslash\widetilde\Lambda}\psi_{\rm per}(w)
 =\sum_{z\in\Lambda}\frac{P_\psi(z)}{e_z}.
\end{equation}
Each lift occurs in exactly $e_z$ cosets, cancelled by its weight;
root-map injectivity ensures that forms are counted once. Also
\[
 \langle P_\psi\rangle=
       \frac1{\vol(\Gamma\backslash\HH)}\int_{\HH}\psi\,d\mu.
\]
Periodisation is essential: there is no restriction $\lambda<1$.
The zero Fourier mode when $\lambda>1$ will contribute a genuine term
$\lambda^2/Y$ to the variance, not an ignorable boundary error.

Finally the Fricke isometry exchanges the two divisor coordinates:
\[
 -\frac1{dw_Q}=\frac{2a+i/\sqrt d}{e}.
\]
After reflection in the imaginary axis this is the point for
$[e,4ad,df]$. Thus one can use the smaller divisor in a suitable cusp.
When $e,f$ are exchanged the original prime functional becomes
$cB-C/d$, rather than $cB-A$; the sieve density must be checked for that
functional too, as in \S\ref{sec:perd}.

%%%%% part B
\section{Sobolev duality}\label{sec:sobolev}
Throughout this section \(\Gamma\) denotes the
\emph{effective} image in \(\PSL_2(\R)\) of a finite-index subgroup of
\(\SL_2(\Z)\), or a conjugate thereof. In particular, kernel sums never
count both signs of a matrix. Write \(V=\vol(\Gamma\backslash\HH)\).
For a \(\Gamma\)-invariant discrete set \(\widetilde\Lambda\) with finitely
many orbits, put
\[
 e_z=|\Gamma_z|,\qquad
 \nu_\Lambda=\sum_{z\in\Gamma\backslash\widetilde\Lambda}e_z^{-1}\delta_z,
 \qquad \#\Lambda=\sum_{z\in\Gamma\backslash\widetilde\Lambda}e_z^{-1}.
\]
These weights are necessary at elliptic points.

\begin{lemma}[Resolvent kernel; standard]\label{L:4.1}
There is a nonnegative continuous radial function $k_0$ on $\HH$, with
\[
 k_0(r)\ll(1+r)e^{-s_0r},\qquad s_0=\frac{1+\sqrt5}{2},
\]
such that the kernel of $(1-\Delta)^{-2}$ on
$L^2(\Gamma\backslash\HH)$ is
$K(z,w)=\sum_{\gamma\in\Gamma}k_0(\operatorname{dist}(z,\gamma w))$.
The sum converges absolutely.
\end{lemma}
\begin{proof}
The Selberg--Harish-Chandra multiplier is
\(h(t)=(5/4+t^2)^{-2}\); see \cite[\S1.8]{Iw} for the transform and its
inversion, and \cite[Theorem 1.14]{Iw} for automorphisation.
It is even, holomorphic in \(|\Im t|<\sqrt5/2\), and integrable against
\(t\tanh(\pi t)\,dt\). Inversion therefore gives a continuous bounded
point-pair kernel, with
$k_0(0)=(4\pi)^{-1}\int_{\R}h(t)t\tanh(\pi t)\,dt$.
For completeness, the decay can also be read directly from the resolvent
kernel. For \(s>1\), the kernel of \((-\Delta+s(s-1))^{-1}\) on \(\HH\) is
\[
 R_s(r)=\frac1{2\pi\sqrt2}\int_r^\infty
       \frac{e^{-(s-1/2)v}}{\sqrt{\cosh v-\cosh r}}\,dv.
\]
Taking the negative derivative in $s(s-1)$, at $s=s_0$, gives
\[
 k_0(r)=\frac1{2\pi\sqrt2(2s_0-1)}
       \int_r^\infty\frac{v e^{-(s_0-1/2)v}}
                            {\sqrt{\cosh v-\cosh r}}\,dv.
\]
This proves positivity. For \(r\ge1\), substitute \(v=r+u\) and use
\(\cosh(r+u)-\cosh r\gg e^r(e^u-1)\); the resulting integral is
$\ll(1+r)e^{-s_0r}$. At \(r=0\) the numerator cancels the
singularity, giving the same boundedness and continuity as inversion.
An orbit of a discrete cofinite group has \(O_{z,w,\Gamma}(e^R)\) points,
with multiplicity, within distance $R$. Since $s_0>1$, the
periodisation is absolutely convergent. The automorphic kernel theorem
identifies it with the multiplier \(h\), which is precisely
\((1-\Delta)^{-2}\).
\end{proof}

\begin{lemma}[Separated sets]\label{L:4.2}
\status{proved}. If distinct points of \(\widetilde\Lambda\) have distance
at least \(r_0>0\), then
\[
 \|\nu_\Lambda\|_{H^{-2}}^2
 :=\langle\nu_\Lambda,(1-\Delta)^{-2}\nu_\Lambda\rangle
 \le C(r_0)\#\Lambda,
\]
where \(C(r_0)\) is independent of the group and the set.
\end{lemma}
\begin{proof}
Using Lemma~\ref{L:4.1} and cancelling the stabiliser multiplicity in
one variable gives
\[
 \langle\nu_\Lambda,(1-\Delta)^{-2}\nu_\Lambda\rangle
 =\sum_{z\in\Gamma\backslash\widetilde\Lambda}\frac1{e_z}
       \sum_{w\in\widetilde\Lambda}k_0(\operatorname{dist}(z,w)).
\]
Indeed each lift of \(w\) occurs \(e_w\) times in the kernel sum.
The balls of radius \(r_0/2\) around distinct lifts are disjoint.
Those with centres within distance \(R\) of \(z\) lie in the ball of
radius \(R+r_0/2\); comparison with its area
\(4\pi\sinh^2((R+r_0/2)/2)\) bounds their number by \(O_{r_0}(e^R)\).
Thus each inner sum is at most
$O_{r_0}(\sum_{m\ge0}(1+m)e^{(1-s_0)m})=O_{r_0}(1)$.
All rearrangements are justified by positivity and this convergent bound.
\end{proof}

\begin{proposition}[Uniform Sobolev duality]\label{L:4.3}
\status{proved}. Suppose \(P\) is smooth and
\(P,\Delta P,\Delta^2P\in L^2(\Gamma\backslash\HH)\). Put
\(\langle P\rangle=V^{-1}\int P\,d\mu\) and
\(P_0=P-\langle P\rangle\). Under Lemma~\ref{L:4.2}'s separation hypothesis,
\[
 \left|\sum_{z\in\Gamma\backslash\widetilde\Lambda}\frac{P(z)}{e_z}
           -\#\Lambda\langle P\rangle\right|
 \le C(r_0)^{1/2}(\#\Lambda)^{1/2}\|(1-\Delta)P_0\|_2.
\]
\end{proposition}
\begin{proof}
Lemma~\ref{L:4.2} says that \((1-\Delta)^{-1}\nu_\Lambda\in L^2\)
with squared norm at most \(C(r_0)\#\Lambda\). Self-adjointness, first
for smooth regularisations of the point masses and then by Sobolev
duality, gives
\[
 \nu_\Lambda(P_0)=
 \langle(1-\Delta)P_0,(1-\Delta)^{-1}\nu_\Lambda\rangle.
\]
Point evaluation agrees with this distributional pairing by local
Sobolev embedding in dimension two. Cauchy--Schwarz proves the claim.
\end{proof}

\begin{corollary}[Heegner forms]\label{L:4.4}
\status{proved}. Let \(\cF\subset\cF_d\) be invariant under a subgroup
\(\Gamma\subset\Gamma_0(d)\) acting on \(w_Q=z_Q/d\), with finitely many
orbits. Proposition~\ref{L:4.3} holds with an absolute constant for its
root set. For $\psi\in C_c^\infty(\HH)$, the periodised Poincar\'e
series of \S\ref{sec:heegner} satisfies its integrability hypotheses,
the counting identity \eqref{A:eq:countidentity}, and
$\langle P_\psi\rangle=V^{-1}\int_{\HH}\psi\,d\mu$.
The same holds after conjugating the group by a real dilation.
\end{corollary}
\begin{proof}
The dilation is a hyperbolic isometry, so Lemma~\ref{L:2.2} supplies the
absolute separation \(r_0=\operatorname{arccosh}(3/2)\).
Unfold the group sum, cancelling elliptic stabilisers as in
Lemma~\ref{L:4.2}, and then unfold translations. The root map is
injective at fixed discriminant, so this counts forms, not just roots.
Unfolding the integral proves the mean formula. The Poincar\'e series
and its Laplacian iterates are smooth and vanish sufficiently high in
every cusp: for a nonstabilising coset the transformed height tends to
zero, below the support of \(\psi\), whereas at infinity the stabilising
coset eventually lies above its support. Thus they satisfy the
integrability assumptions of Proposition~\ref{L:4.3}.
\end{proof}

\section{The spectral large sieve and the variance of box Poincar\'e series}
\label{sec:sieve}
Use the effective group $\Gamma_{M,q}$ of \S\ref{sec:heegner},
$M=4dq^2$, with squarefree $(q,2d)=1$ and cusp width one. Write
$h_q=[\Gamma_0(M):\Gamma_{M,q}]$: $h_1=1$ and
$h_q=\varphi(q)/2$ for $q>1$. Its $h_q$ character spaces are the
even nebentypus spaces modulo $q$.

\subsection*{The imported spectral estimates}
We specify normalisations to distinguish a cusp width from the parameter
in the large sieve. At a singular cusp \(\mathfrak a\), choose a scaling
matrix giving period one, and write the expansion of an orthonormal
weight-zero Maass form as
\begin{equation}\label{B:eq:fourier}
 u_j(\sigma_{\mathfrak a}z)
 =\sqrt y\sum_{n\ne0}\rho_{j\mathfrak a}(n)
                 K_{it_j}(2\pi|n|y)e^{2\pi inx}.
\end{equation}
This is DI's normalisation \cite[(1.34), p.~231]{DI}.
The Eisenstein coefficients \(\rho_{\mathfrak c\mathfrak a}(n,t)\)
will use the identical \(\sqrt yK_{it}\) convention. The DI cusp
parameter for \(\Gamma_0(M)\) is
\[
 \mu(\mathfrak a)=\frac{(v,M/v)}M
 \quad\text{if }\mathfrak a\sim u/v,\ v\mid M,\ (u,v)=1;
 \qquad \mu(\infty)=M^{-1}.
\]
It is \emph{not} the literal translation width at infinity, which is one.
For \(K\ge1\), \(H\ge1/2\), and \(a_n\) supported on \(H<n\le2H\),
DI \cite[Theorem 2, (1.29)--(1.30), p.~230]{DI} bounds the cuspidal
expression
\[
 \sum_{|t_j|\le K}\frac1{\cosh(\pi t_j)}
       \left|\sum_n a_n\rho_{j\mathfrak a}(n)\right|^2
\]
and its continuous-spectrum analogue by
\((K^2+\mu(\mathfrak a)H^{1+\varepsilon})\sum_n|a_n|^2\), for trivial
nebentypus. With an even character of conductor \(q_0\mid q\), the bound is
(Drappeau's \(q_0\) is the modulus of \(\chi\); on \(\Gamma_0(M)\) the
multiplier \(\chi(\gamma_{22})\) depends only on the primitive character
inducing \(\chi\), so \(q_0\) may be taken to be the conductor; R114 repair)
\begin{equation}\label{B:eq:large-sieve}
 \ll_\varepsilon
 (K^2+q_0^{1/2}\mu(\mathfrak a)H^{1+\varepsilon})\sum_n|a_n|^2.
\end{equation}
Here the continuous analogue means
\(\sum_{\mathfrak c}(4\pi)^{-1}\int_{-K}^K
\cosh(\pi t)^{-1}|\sum_n a_n\rho_{\mathfrak c\mathfrak a}(n,t)|^2dt\),
summed over all singular cusps. This is
\cite[Proposition 4.7, (4.24)--(4.25)]{Dr}, in the arXiv numbering.
Drappeau uses Whittaker coefficients: since
\(W_{0,it}(4\pi|n|y)=2\sqrt{|n|y}K_{it}(2\pi|n|y)\),
\(\rho^{\rm DI}(n)=2\sqrt{|n|}\rho^{\rm Dr}(n)\).
Thus his \(\sqrt n\rho^{\rm Dr}(\pm n)\) is exactly our coefficient
up to the constant 2; his factor \((1+|t|)^{\pm\kappa}\) is one because
\(\kappa=0\). These estimates hold separately for either sign of \(n\).

For later comparison, DI \cite[Theorem 5, (1.38), p.~232]{DI} states,
for trivial character, \(X\ge1\), and exceptional parameters
\(t_j=i\sigma_j\), \(0<\sigma_j<1/2\),
\begin{equation}\label{B:eq:exceptional-sieve}
 \sum_{j\,\mathrm{exc}}X^{2\sigma_j}
   \left|\sum_{H<n\le2H}a_n\rho_{j\mathfrak a}(n)\right|^2
 \ll_\varepsilon
 (1+\sqrt{\mu(\mathfrak a)HX})
 (1+\sqrt{\mu(\mathfrak a)H^{1+\varepsilon}})\|a\|_2^2.
\end{equation}
The nebentypus analogue is \cite[Lemma 4.8]{Dr}, with second factor
\(1+(q_0\mu(\mathfrak a)H)^{1/2+\varepsilon}\).
The displayed DI statements have been checked against the scan; the
Drappeau numbering and coefficient conversion refer to arXiv:1504.05549,
not to an unchecked published numbering.

\subsection*{The box variance}
Take \(\phi\in C_c^\infty([-2,-1])\), \(W\in C_c^\infty([1,2])\), and
\[
 \psi(u)=\phi(x/\lambda)W(y/Y),\qquad
 L_* =\log(2+(\lambda Y)^{-1}+Y^{-1}+M).
\]
Assume \(\lambda>0\) and \(Y\le\lambda L_*^{-3}\). Periodisation is always
understood; no restriction \(\lambda<1\) is imposed. For nonnegative
nonzero weights its area is \(\asymp\lambda/Y\). In applications
\(\lambda\asymp A/(qF)\), \(Y\asymp1/(qF\sqrt d)\), and \(L_*\ll\cL\).

\begin{proposition}[Box variance]\label{L:5.1}
\status{conditional} on (SEL). On \(\Gamma_{M,q}\backslash\HH\),
\begin{equation}\label{B:eq:variance}
 \|(1-\Delta)(P_\psi-\langle P_\psi\rangle)\|_2^2
 \ll_{\phi,W,\varepsilon}L_*^C
 \left\{\frac\lambda Y
   \left(1+\frac{q^{1/2}}M(L_*/\lambda)^{1+\varepsilon}\right)
       +\frac{\lambda^2}Y\right\}.
\end{equation}
The constant is uniform in \(d,q,\lambda,Y\), and depends on only finitely
many seminorms of the two fixed weights.
\end{proposition}
\begin{proof}
\emph{Reduction and spectral expansion.}
Differentiation commutes with periodisation and
\[
 (1-\Delta)\psi
 =\phi(x/\lambda)(W(v)-v^2W''(v))
       -(Y/\lambda)^2\phi''(x/\lambda)v^2W(v),\qquad v=y/Y.
\]
Moreover \(\int\Delta\psi\,d\mu=0\). Since \(Y/\lambda\le1\), it suffices
to bound \(\|P_\psi-\langle P_\psi\rangle\|_2^2\) by the right side of
\eqref{B:eq:variance}, for arbitrary fixed smooth weights of this type.
For each even \(\chi\bmod q\), let
\[
 P_\chi=\sum_{\gamma\in\Gamma_\infty\backslash\Gamma_0(M)}
            \overline{\chi(\gamma)}\psi_{\rm per}(\gamma u).
\]
Character orthogonality gives \(P_\psi=h_q^{-1}\sum_\chi P_\chi\) and
\begin{equation}\label{B:eq:char-norm}
 \|P_\psi-\langle P_\psi\rangle\|_{\Gamma_{M,q}}^2
 =\frac1{h_q}\sum_\chi
       \|P_\chi-\ind_{\chi=1}\langle P_1\rangle\|_{\Gamma_0(M),\chi}^2.
\end{equation}
The spectral expansion, with Eisenstein measure \(dt/(4\pi)\), is the
usual Parseval formula \cite[Theorem 7.3]{Iw} (with the character
spaces and singular cusps as in \cite[\S4.1]{Dr}).
\emph{This is the sole use of (SEL):} after removing the constant,
all discrete parameters are real, and there is no nonconstant residual
spectrum. Thus the squared norm for a fixed character is the sum of
\(|\langle P_\chi,u_j\rangle|^2\) over cusp forms and
\((4\pi)^{-1}\sum_{\mathfrak c}\int_\R
|\langle P_\chi,E_{\mathfrak c}(\cdot,1/2+it)\rangle|^2dt\).

\emph{Unfolding and Mellin inversion.}
Write \(\widehat\phi(\xi)=\int_\R\phi(v)e^{-2\pi i\xi v}dv\).
The periodised box has Fourier coefficient
\(\lambda\widehat\phi(\lambda n)W(y/Y)\), even when \(\lambda>1\).
Using \eqref{B:eq:fourier}, unfolding gives
\[
 \langle P_\chi,u_j\rangle
 =\sum_{n\ne0}\overline{\rho_j(n)}\lambda\widehat\phi(\lambda n)I_{t_j}(n),
 \quad I_t(n)=\int_0^\infty W(y/Y)y^{-3/2}K_{it}(2\pi|n|y)\,dy.
\]
Set
\[
 G_t(s)=\Gamma((s+it)/2)\Gamma((s-it)/2),\qquad
 \mathcal W(s)=\int_0^\infty W(v)v^{-3/2-s}dv.
\]
Write also
\[
 B_j(s)=\sum_{n\ne0}\overline{\rho_j(n)}
              \lambda\widehat\phi(\lambda n)|n|^{-s}.
\]
The function \(\mathcal W\) is entire and rapidly decreasing vertically,
uniformly on fixed strips. Mellin inversion of the Bessel function gives,
for \(\sigma>0\),
\begin{equation}\label{B:eq:mellin}
 \langle P_\chi,u_j\rangle
 =\frac{Y^{-1/2}}{8\pi i}\int_{(\sigma)}
        G_{t_j}(s)\mathcal W(s)(\pi Y)^{-s}B_j(s)\,ds.
\end{equation}
Absolute convergence follows from the rapid decay of
\(\widehat\phi\), the polynomial growth of Fourier coefficients, and
Stirling's formula.

\emph{Uniform spectral bounds.}
Put \(\sigma_0=1/L_*\). Stirling's formula, including the simple poles
of \(\Gamma\), gives
\begin{align*}
 |G_t(\sigma_0+iv)|&\ll L_*^2e^{-\pi|v|/2} &&(|t|\le1),\\
 |G_t(-1+iv)|&\ll e^{-\pi|t|/2}(1+|t|)^{-2}(1+|v|)^4
                                      &&(|t|\ge1).
\end{align*}
For the second bound, if \(|v|\le|t|/2\) both gamma arguments have
imaginary part comparable to \(|t|\); otherwise
\(1+|t|\ll1+|v|\) and the exponential is
\(e^{-\pi\max(|v|,|t|)/2}\). On this line both real parts are \(-1/2\),
a fixed distance from gamma poles. Since \(Y^{-\sigma_0}\ll1\),
Cauchy--Schwarz in \eqref{B:eq:mellin} yields, for \(|t_j|\le1\),
\[
 |\langle P_\chi,u_j\rangle|^2
 \ll Y^{-1}L_*^4\int_\R |\mathcal W(\sigma_0+iv)|
                                    |B_j(\sigma_0+iv)|^2\,dv.
\]
For \(|t_j|\ge1\), shift to \(\Re s=-1\), crossing only \(s=\pm it_j\).
The residue contribution is
\[
 \frac{Y^{-1/2}}2\sum_\pm
       \Gamma(\pm it_j)\mathcal W(\pm it_j)(\pi Y)^{\mp it_j}B_j(\pm it_j).
\]
Let \(\widetilde B_j(s)=(\pi Y)^{-s}B_j(s)\).
Using \(|\Gamma(it)|^2=\pi/(|t|\sinh(\pi|t|))\), for any fixed large
\(B\) the squared coefficient is at most
\begin{align}\label{B:eq:high-t}
 \frac{C_BY^{-1}}{\cosh(\pi t_j)}\bigg\{&
 (1+|t_j|)^{-B}\sum_\pm|B_j(\pm it_j)|^2\notag\\
 &+(1+|t_j|)^{-4}\int_\R
  |\mathcal W(-1+iv)|(1+|v|)^4|\widetilde B_j(-1+iv)|^2\,dv\bigg\}.
\end{align}

\emph{The large sieve, including both tails.}
Split \(n\) by sign and into dyadic blocks \(H<|n|\le2H\),
\(H=2^{i-1}\), \(i\ge0\). With
\(w_H=1+|\log(\lambda H)|\), one has \(\sum_Hw_H^{-2}\ll1\).
Cauchy--Schwarz across blocks and \eqref{B:eq:large-sieve} imply
\[
 \mathcal S(K;b):=\sum_{|t_j|\le K}\frac{|\sum_{n\ne0}\overline{\rho_j(n)}b_n|^2}
                                             {\cosh(\pi t_j)}
 \ll_\varepsilon\sum_H w_H^2
       (K^2+q^{1/2}M^{-1}H^{1+\varepsilon})\|b^{(H)}\|_2^2.
\]
Complex conjugation of the coefficient vector causes no change in this
inequality; either sign is covered by Drappeau's statement (or reflection
and replacement of \(\chi\) by \(\bar\chi\)). For
\(b_n(s)=\lambda\widehat\phi(\lambda n)|n|^{-s}\), at
\(\Re s=\sigma_0\) or \(0\), the block norm is
\(\ll_B\lambda^2H(1+\lambda H)^{-B}\).
Summing the geometric tails gives, with
\(Z=q^{1/2}M^{-1}\lambda^{-1-\varepsilon}\),
\begin{equation}\label{B:eq:block-sieve}
 \mathcal S(K;b(s))\ll L_*^2\lambda(K^2+Z).
\end{equation}
At \(\Re s=-1\), the coefficients of \(\widetilde B\) have modulus
\(\pi Y|n|\lambda|\widehat\phi(\lambda n)|\); their bound is
\eqref{B:eq:block-sieve} multiplied by \((Y/\lambda)^2\le1\), after
increasing the decay exponent. This also proves the claims when
\(\lambda>1\): then every nonzero mode lies in the rapidly decreasing tail.

For the varying vectors \(b_n(\pm it_j)\), on each unit interval
\([k,k+1]\) use the elementary inequality
\[
 |B_j(\pm it)|^2\le2|B_j(\pm ik)|^2+
              2\int_k^{k+1}|\partial_\tau B_j(\pm i\tau)|^2d\tau.
\]
Differentiating multiplies a coefficient by \(\log|n|\), bounded on a
block by \(O(L_*+|\log(\lambda H)|)\); the extra factors are absorbed
in the same block sums. Consequently the sum on \(K<|t_j|\le2K\)
for these vectors is \(\ll K L_*^4\lambda(K^2+Z)\).
Insert these bounds in \eqref{B:eq:high-t} and the small-\(t\) estimate.
Dyadic summation over \(K\ge1\) converges, since the integral term has
weight \(K^{-4}\) and the residue term has arbitrary weight \(K^{-B}\).
Vertical integrals converge by the decay of \(\mathcal W\). We obtain
\[
 \sum_j|\langle P_\chi,u_j\rangle|^2
       \ll L_*^C\frac\lambda Y(1+Z).
\]
No truncation at logarithmic frequency or height was made.

\emph{The continuous spectrum and its zero mode.}
For nonzero modes the identical argument uses the continuous version of
\eqref{B:eq:large-sieve}. In the unit-interval argument freeze the
Eisenstein parameter in the Fourier coefficients and differentiate only
the auxiliary variable in \(|n|^{-i\tau}\), then evaluate on the diagonal
\(\tau=t\); no derivative of an Eisenstein coefficient is required.
The constant term at infinity is
\(\delta_{\mathfrak c\infty}y^{1/2+it}
 +\Phi_{\mathfrak c\infty}(1/2+it)y^{1/2-it}\).
Its pairing with the box is
\[
 \lambda\widehat\phi(0)Y^{-1/2}
 \{\delta_{\mathfrak c\infty}Y^{-it}\mathcal W(it)
  +\overline{\Phi_{\mathfrak c\infty}(1/2+it)}Y^{it}\mathcal W(-it)\}.
\]
Unitarity of the scattering matrix on the unitary axis gives
\(\sum_{\mathfrak c}|\Phi_{\mathfrak c\infty}(1/2+it)|^2=1\).
Squaring, integrating and summing therefore costs \(O(\lambda^2/Y)\),
independently of the number of cusps. The cross term with nonzero modes
is absorbed by \(|a+b|^2\le2|a|^2+2|b|^2\).
Finally \eqref{B:eq:char-norm} averages exactly \(h_q\) character bounds;
there is no loss from their number. Enlarging
\(\lambda^{-1-\varepsilon}\) to \((L_*/\lambda)^{1+\varepsilon}\) proves
\eqref{B:eq:variance}.
\end{proof}

Without (SEL), the terms with \(t_j=i\sigma_j\) remain. On a block
\(|n|\asymp H\), the Bessel factor has additional size
\((HY)^{-\sigma_j}\). Thus \eqref{B:eq:exceptional-sieve} is used with
\(X=(HY)^{-1}\) when \(HY\le1\), not with \(X=Y^{-1}\).
It does not recover Proposition~\ref{L:5.1} in the required cells;
Section~\ref{sec:uncond} records precisely the resulting unconditional
limitation. Neither the Sobolev argument nor the regular-spectrum large
sieve depends on (SEL).

\section{Per-\texorpdfstring{$d$}{d} counts}\label{sec:perd}
Fix $c\ge1$ and use $\cF_d^I$, $\Gamma'_q$ from \S\ref{sec:heegner},
with $n(Q)=cB-A=4acd-f$. In $[A,B,C]$, $A$ is a form coefficient,
not the dyadic parameter; the parity and sieve invariance were proved there.

For a prime \(\ell\nmid2d\), put \(\chi_d(\ell)=(-d/\ell)\) (this agrees
with \((\frac{-4d}{\ell})\) of Lemma~\ref{L:8.3}; R114 repair) and define
\(g_{c,d}(\ell)\) as the proportion on the quadric
\(B^2-4AC=-4d\) satisfying \(cB-A=0\). Directly,
\begin{equation}\label{B:eq:density}
 g_{c,d}(\ell)=
 \begin{cases}
 (\ell-1)/(\ell^2+\chi_d(\ell)\ell),&\ell\nmid c,\\
 (1+\chi_d(\ell))/(\ell+\chi_d(\ell)),&\ell\mid c.
 \end{cases}
\end{equation}
Extend multiplicatively to squarefree \(q\) coprime to \(2d\).
The denominator counts the quadric: for \(A\ne0\) there are
\(\ell(\ell-1)\) choices, while \(A=0\) gives
\(\ell(1+\chi_d(\ell))\). If \(\ell\nmid c\), the imposed equation
forces \(B\ne0\), giving exactly \(\ell-1\) choices. If \(\ell\mid c\),
it forces \(A=0\), giving the second numerator. Primes dividing \(2d\)
never divide \(n\): both \(f\) and \(n\) are odd, and \(f\) is a unit
at every prime dividing \(d\).

\begin{lemma}[Factorisation of the orbital main term]\label{L:6.1}
\status{proved}. Let \(\#\Lambda_d(q)\) be the orbifold count of
\(\{Q\in\cF_d^I:q\mid n(Q)\}\) modulo \(\Gamma'_q\), and put
\(V(q)=\vol(\Gamma'_q\backslash\HH)\). For squarefree \((q,2d)=1\),
\[
 \frac{\#\Lambda_d(q)}{V(q)}
 =g_{c,d}(q)\frac{\#\Lambda_d(1)}{V(1)}.
\]
\end{lemma}
\begin{proof}
There is nothing to prove for \(q=1\). Reduction of \(\Gamma'_1\)
onto \(\SL_2(\Z/q\Z)\) is surjective: it contains \(\Gamma(2d)\),
and the assertion follows from CRT and surjectivity of reduction for
\(\SL_2(\Z)\). The kernel is \(\Gamma'_q\). Thus the index of effective
groups is
\(I_q=|\SL_2(\Z/q\Z)|/2\), since \(-I\) belongs to \(\Gamma'_1\)
but not to \(\Gamma'_q\).
At a prime \(\ell\mid q\), the stabiliser of a nondegenerate binary form
in \(\SL_2(\mathbb F_\ell)\) is its special orthogonal torus, of order
\(\ell-\chi_d(\ell)\): it is either the split torus or the norm-one
subgroup of \(\mathbb F_{\ell^2}^{\times}\).
Its orbit consequently has \(\ell^2+\chi_d(\ell)\ell\) elements, the
size of the entire quadric. The fraction satisfying the sieve condition
is therefore \(g_{c,d}(\ell)\), independently of the original orbit.
CRT gives the product of these fractions.
For an orbit with stabiliser order \(e_Q\), the weighted number of
sieved \(\Gamma'_q\)-orbits above it is
\(e_Q^{-1}I_qg_{c,d}(q)\), by orbit--stabiliser (equivalently, sum over
the cosets before dividing by the stabiliser). Summing over the original
orbits, and using \(V(q)=I_qV(1)\), proves the formula.
\end{proof}

Choose a fixed physical box, so that in the coordinates \(u=w/(2q)\)
its weights are \(\psi_q(u)=\psi_1(qu)\). Accordingly
\(\lambda_q=\lambda_1/q\), \(Y_q=Y_1/q\), and
\(\int\psi_q\,d\mu=\int\psi_1\,d\mu\). Define
\[
 \mathfrak M_d=\frac{\#\Lambda_d(1)}{V(1)}\int_{\HH}\psi_q\,d\mu,
 \qquad \kappa_c(q)=2^{\omega((q,c))}.
\]
The first quantity is independent of the sieve modulus. Here and below
\(\omega(m)\) counts the distinct prime factors of \(m\).

\begin{theorem}[Per-\(d\) count]\label{L:6.2}
\status{conditional} on (SEL). Let \(q\) be squarefree, \((q,2d)=1\),
and let \(\psi_q\) be a box as in Proposition~\ref{L:5.1}, with
\(\lambda_q\asymp A/(qF)\), \(Y_q\asymp1/(qF\sqrt d)\), \(A,F\ge1\).
Assume explicitly $Y_q\le\lambda_q L_*^{-3}$ and polynomial parameter
ranges in $N$, so $L_*\ll\cL$.
Then, for \(0<\varepsilon<1/2\),
\begin{align}\label{B:eq:per-d}
 &\left|\sum_{\substack{Q\in\cF_d^I\\q\mid n(Q)}}\psi_q(u_Q)
             -g_{c,d}(q)\mathfrak M_d\right|\notag\\
 &\quad\ll_\varepsilon \cL^C q\kappa_c(q)^{1/2}(\#\Lambda_d(1))^{1/2}
   \left\{A\sqrt d\left(1+\frac A{qF}\right)
                     +\frac{F^{1+\varepsilon}}{\sqrt d}\right\}^{1/2}.
\end{align}
In particular the factor \(\kappa_c(q)^{1/2}\) is absent when \((q,c)=1\),
as in the sieve application. There is no lower bound on \(F/A\).
\end{theorem}
\begin{proof}
Corollary~\ref{L:4.4} applies to the sieved subset, because separation
is preserved by passage to a subset and by the dilation to \(u\).
Lemma~\ref{L:6.1} identifies its mean as \(g_{c,d}(q)\mathfrak M_d\).
In Proposition~\ref{L:5.1},
\[
 \frac\lambda Y\asymp A\sqrt d,\qquad
 \frac{\lambda^2}Y\asymp A\sqrt d\frac A{qF},\qquad
 \frac{q^{1/2}}{MY}\lambda^{-\varepsilon}
 \ll\frac{F^{1+\varepsilon}}{\sqrt d}
                 q^{-1/2+\varepsilon}A^{-\varepsilon}
 \le\frac{F^{1+\varepsilon}}{\sqrt d}.
\]
The remaining factor from Sobolev duality is \((\#\Lambda_d(q))^{1/2}\).
For \(\ell\nmid c\),
\[
 |\SL_2(\mathbb F_\ell)|g_{c,d}(\ell)
       =(\ell-\chi_d(\ell))(\ell-1)\le\ell^2;
\]
for \(\ell\mid c\) it is either zero or \(2\ell(\ell-1)\le2\ell^2\).
The exact identity from Lemma~\ref{L:6.1} therefore gives
\(\#\Lambda_d(q)\le\kappa_c(q)q^2\#\Lambda_d(1)\), also for \(q=1\).
Taking square roots proves \eqref{B:eq:per-d}.
\end{proof}

The term \(A/(qF)\) is the cost of wrapping around several periods;
it carries no \(N^\varepsilon\) loss. One may instead use the \(e\)-cusp:
\(w\mapsto-1/(dw)\) sends the point to
\((2a+i/\sqrt d)/e\), and reflection gives the same boxes.
Interchanging \(e,f\) replaces the sieve equation \(A=cB\) by
\(C/d=cB\). The involution
\((A,B,C)\mapsto(C/d,B,dA)\) is a bijection of the quadric modulo every
prime not dividing \(2d\); hence the densities and the estimate are
unchanged.

\begin{lemma}[Class-number bound]\label{L:6.3}
\status{proved}. With \(h(-4d)\) the number of primitive proper classes
of discriminant \(-4d\), and
\(r(d)=\prod_{p^k\parallel d}p^{\lfloor k/2\rfloor}\),
\[
 \#\Lambda_d(1)\le6h(-4d)r(d),\qquad
 \sum_{d\le D}\#\Lambda_d(1)\ll D^{3/2}\log^2(2D).
\]
\end{lemma}
\begin{proof}
Every Type~I form is primitive: a prime dividing \(f\) divides neither
\(2ad\) nor \(4ad\), by \(ef-4a^2d=1\).
For a primitive class represented by \(Q_0\), its admissible level-\(d\)
cosets are bounded by the number of
\((t:s)\in\mathbb P^1(\Z/d\Z)\) with \(Q_0(t,s)\equiv0\pmod d\).
Indeed, the second column of a matrix representing a coset modulo
\(\Gamma^0(d)\) gives this projective point; the condition that the
transformed last coefficient be divisible by \(d\) is necessary.
Passing from \(\Gamma_0(d)\) to \(\Gamma'_1\) costs index at most six
(exactly six for odd \(d\), four for even \(d\)); imposing Type~I parity
can only decrease the count. Stabiliser weights only improve this upper bound.

For \(p^k\parallel d\), make an invertible change of variables so that
\(Q_0=[A_0,B_0,C_0]\) has \(p\nmid A_0\). This is possible by primitivity;
at 2, \(B_0\) is even and one of the outer coefficients is odd.
Points with \(p\mid s\) cannot be zeros. At an odd prime, zeros
\((t:1)\) satisfy
\((2A_0t+B_0)^2\equiv0\pmod{p^k}\), since the discriminant is \(-4d\).
There are exactly \(p^{\lfloor k/2\rfloor}\) possibilities.
At 2, write \(B_0=2B'_0\); the identity
\[
 A_0Q_0(t,1)=(A_0t+B'_0)^2+d
\]
shows that the condition is
\(A_0t+B'_0\equiv0\pmod{2^{\lceil k/2\rceil}}\), again giving
\(2^{\lfloor k/2\rfloor}\) possibilities. CRT proves the first assertion.
The classical imaginary-quadratic class-number formula
\cite{IK} is
\(h(D_0)=w(D_0)\sqrt{|D_0|}L(1,\chi_{D_0})/(2\pi)\),
where \(w(D_0)\le6\) and \(\chi_{D_0}\) is the Kronecker character,
possibly imprimitive. Its nonprincipal character sums are bounded by
\(|D_0|\); splitting the Dirichlet series there and using partial
summation gives \(L(1,\chi_{D_0})\ll\log(2|D_0|)\).
Thus \(h(-4d)\ll\sqrt d\log(2d)\), including nonfundamental discriminants.
Finally,
\[
 \sum_{d\le D}r(d)
 \le\sum_{m\le\sqrt D}m\,\#\{d\le D:m^2\mid d\}
 \le D\sum_{m\le\sqrt D}\frac1m\ll D\log(2D).
\]
Multiplying by \(\sqrt D\log(2D)\) proves the second assertion.
\end{proof}

There is no uniform per-\(d\) relative-error assertion: the orbital
main term contains a possibly small class number. Section~\ref{sec:assembly}
uses only absolute errors and
\(\sum_{d\asymp D}\sqrt{\#\Lambda_d(1)}\ll D^{5/4}\log(2D)\),
by Cauchy--Schwarz and Lemma~\ref{L:6.3}. Relative to the box mass
\(AD\), the first summed error in \eqref{B:eq:per-d} has size
\(\cL^Cq(D/A)^{1/2}(1+A/(qF))^{1/2}\) when \((q,c)=1\).
This is an averaged absolute-error calculation, not a lower bound for
an individual \(\mathfrak M_d\).

%%%%% part C
\section{The per-\texorpdfstring{$a$}{a} Weil count}\label{sec:pera}
The complementary estimate is unconditional.  It counts a modular hyperbola
rather than a packet of forms.  Write $e(t)=\exp(2\pi it)$ in this section,
and put
\[
 S(h,k;m)=\sum_{x\bmod m}^{*}e\bigl((hx+k\bar x)/m\bigr).
\]
All smooth weights below have fixed derivative bounds.  For upper bounds
we take them nonnegative, with positive integral.

\begin{proposition}[the count $(K_a)$; \status{proved}]\label{L:7.1}
Let $a,c\geq1$, let $q$ be squarefree with $(q,2a)=1$, and set
$m=4a^2$, $r=mq$.  For $D,E\geq1$ put $F=mD/E$, and let
$W_1,W_2\in C_c^\infty((1,2))$.  Define
\[
 S_a(q)=\sum_{\substack{e,f\geq1\\ef\equiv1\ (m)\\q\mid4acd-f}}
 W_1(e/E)W_2(d/D),\qquad d=(ef-1)/m.
\]
Then $S_a(q)=g'_{c,a}(q)X_a+R_a(q)$, where
\begin{align}
 X_a&=\frac{\varphi(m)}{m^2}\iint_{\R^2}
 W_1(x/E)W_2\bigl((xy-1)/(mD)\bigr)\,dx\,dy
 \asymp D\frac{\varphi(m)}m,\label{C:eq:Ka-main}\\
 |R_a(q)|&\ll 2^{\omega(q)}\tau(m)^2
 \left(qa+\frac DE+\frac DF\right),\label{C:eq:Ka-error}
\end{align}
and $g'_{c,a}$ is multiplicative on squarefree numbers, with
\[
 g'_{c,a}(\ell)=
 \begin{cases}(\ell-1)/\ell^2,&\ell\nmid c,\\1/\ell,&\ell\mid c.
 \end{cases}
\]
The constants depend only on the weights.
\end{proposition}
\begin{proof}
The identity
\[
 an=f(ce-a)-c
\]
shows that the residue conditions on $(e,f)$ modulo $r$ are the CRT
product of
\[
 \cS_m=\{(e,f):ef=1\pmod m\},\qquad
 \cS_q=\{(e,f):f(ce-a)=c\pmod q\}.
\]
Put $W(x,y)=W_1(x/E)W_2((xy-1)/(mD))$.  On its support
$x\asymp E$, $y\asymp F$, and
$\partial_x^i\partial_y^jW\ll_{i,j}E^{-i}F^{-j}$.
Two-dimensional Poisson summation gives
\begin{equation}\label{C:eq:poisson}
 \begin{aligned}
 S_a(q)&=r^{-2}\sum_{h,k\in\Z}\widehat W(h/r,k/r)\mathcal K(h,k),\\
 |\widehat W(h/r,k/r)|&\ll_B EF
 (1+E|h|/r)^{-B}(1+F|k|/r)^{-B}.
 \end{aligned}
\end{equation}
The sign in the finite Fourier transform is immaterial; with the
positive sign, CRT gives
\[
 \mathcal K(h,k)=S(\bar qh,\bar qk;m)T_q(\bar mh,\bar mk),
 \quad T_q(u,v)=\sum_{(e,f)\in\cS_q}e((ue+vf)/q).
\]
The prime factors of $T_q$ have CRT-rescaled arguments.  For a prime
$\ell\mid q$, if $\ell\nmid c$, the substitution $w=ce-a$ gives
\[
 T_\ell(u,v)=e(u\bar ca/\ell)S(u\bar c,vc;\ell).
\]
If $\ell\mid c$, then $f=0$ and
$T_\ell(u,v)=\ell\ind_{\ell\mid u}$.  In particular
$|T_q(u,v)|\leq2^{\omega(q)}q$ (the trivial cardinality bound suffices).
At $(h,k)=(0,0)$ we have $|\cS_m|=\varphi(m)$ and
$|\cS_\ell|=\ell-1$ or $\ell$, respectively.  This gives exactly
\eqref{C:eq:Ka-main} times $g'_{c,a}(q)$.

We retain the axial frequencies: they are important if a divisor is
small.  For $B>2$ and every integer $v\geq1$,
\begin{equation}\label{C:eq:frequency-tail}
 \sum_{\substack{k\ne0\\v\mid k}}(1+F|k|/r)^{-B}\ll_B\frac r{vF}.
\end{equation}
Indeed, if $vF\leq r$, compare with an integral; otherwise the left side
is $O_B((r/(vF))^B)$.  Since
$|S(0,\bar qk;m)|=|c_m(k)|\leq(k,m)$, expanding
$(k,m)\leq\sum_{v\mid(k,m)}v$ in \eqref{C:eq:frequency-tail} gives
\[
 \sum_{k\ne0}(k,m)(1+F|k|/r)^{-B}\ll\tau(m)r/F.
\]
Thus the frequencies $h=0\ne k$ contribute at most
$2^{\omega(q)}\tau(m)E/m=2^{\omega(q)}\tau(m)D/F$;
$k=0\ne h$ contributes at most $2^{\omega(q)}\tau(m)D/E$.

For $hk\ne0$ use the classical Weil--prime-power bound
\cite{Weil}, in the convenient form
\[
 |S(u,v;m)|\leq3\tau(m)m^{1/2}(u,v,m)^{1/2}.
\]
The constant allows the prime $2$.  Expanding the gcd as before yields
\begin{align*}
 &\sum_{h,k\ne0}(h,k,m)^{1/2}
 (1+E|h|/r)^{-B}(1+F|k|/r)^{-B}\\
 &\hspace{20mm}\leq
 \sum_{v\mid m}v^{1/2}\,O_B\left(\frac r{vE}\frac r{vF}\right)
 \ll_B\frac{r^2}{EF},
\end{align*}
since $\sum_{v\geq1}v^{-3/2}<\infty$.  This remains valid when
$vE>r$ or $vF>r$.  Inserting it in \eqref{C:eq:poisson} bounds the
nonaxial contribution by
$O(2^{\omega(q)}q\tau(m)m^{1/2})$.  As $m^{1/2}=2a$, the claimed
error follows, even with one power of $\tau(m)$ here.
\end{proof}

For later use, if $\min(E,F)\geq N^{c_0}$, summation over $a\asymp A$
and squarefree $q\leq Q$ gives
\begin{equation}\label{C:eq:Ka-summed}
 \sum_{a\asymp A}\sum_{\substack{q\leq Q\\(q,2a)=1}}
 \mu^2(q)3^{\omega(q)}|R_a(q)|
 \ll\cL^C\bigl(Q^2A^2+QADN^{-c_0}\bigr).
\end{equation}
Here we used the elementary divisor-moment bound
$\sum_{a\asymp A}\tau(4a^2)^2\ll A(\log(2A))^C$.
Measured against the scale $AD$, the error is
$\cL^C(Q^2A/D+QN^{-c_0})$.  This gives a useful sieve level on the
side $D>A$.  The axial errors explain why a uniform density assertion
with no lower bound on $E,F$ would be false.  No small-divisor substitute
for this proposition is used below.

\section{Assembly}\label{sec:assembly}
We begin with the elementary accounting fact that permits a sieve level
to approach $1$ at the boundary $A=D$.

\begin{lemma}[linearly degenerating savings; \status{proved}]\label{L:1.1}
Let $L\geq2$, and let $m_{jk}\geq0$, $1\leq j,k\leq\lfloor L\rfloor$,
with $\sum_km_{jk}\leq\mathcal M$ for each $j$.  Then
\[
 \sum_{j,k}m_{jk}\min(1/j,1/k)\leq\mathcal M(1+\log L).
\]
If also $m_{jk}\leq\mathcal M/L$, the same sum is at most $2\mathcal M$.
(R114 repair: $\mathcal M$ is a mass bound, not the level $M=4dq^2$.)
\end{lemma}
\begin{proof}
The first assertion follows by using $1/j$ and summing the harmonic
series.  For the second, group pairs by $h=\max(j,k)$:
\[
 \sum_{j,k\leq L}\min(1/j,1/k)
 =\sum_{h\leq L}\frac{2h-1}{h}\leq2L.
\]
Multiply by $\mathcal M/L$.  Fixed constant multiples of either saving only
change the resulting absolute constant.
\end{proof}
Below, $j$ indexes a $c$-block and $k$ its distance from $A=D$.
Layer masses are $O(N\cL)$, so take $\mathcal M\asymp N\cL^2$ and $L\asymp\cL$.
The block $j=0$ is treated separately, never with $1/j$.

\begin{theorem}[\status{conditional} on (SEL)]\label{L:8.1}
Assume \textup{(SEL)} as stated in Section~\ref{sec:intro}, uniformly for
$\Gamma_0(4dq^2)$ and the even nebentypus characters modulo $2q$.
Then, relative to the cited inputs listed in Section~\ref{sec:intro}
(R114 repair),
\[
 \sum_{p\leq N}f_I(p)\ll N\log^2N.
\]
\end{theorem}
\begin{proof}
We work with $N/2<p\leq N$ and write $\cL=\log N$.
The reduction of Section~\ref{sec:reduction}, from
\cite[Proposition~2.2 and Lemma~2.8]{ET}, gives
$f_I(p)\leq2\sum_cw_c(p)$.  Fix small constants
$0<\eta\leq\eta_1$; all subsequent constants may depend on them.

\smallskip\noindent
\emph{1. Removing the short-progression ranges.}
Proposition~\ref{M:3.8} bounds the part $c>N^\eta$ by $O_\eta(N\cL^2)$.
For $c\leq N^\eta$, Lemma~\ref{L:8.2}
disposes, all $c$ at once, of tuples with
\[
 \min(4ab,4acf,4cdf)\leq N^{1-\eta_1}.
\]
Indeed $bf=na+c$ and $f\leq2n$ imply $b\geq a/2$, so $4ad$ and
$4bd$ are already $\gg N^{1-\eta}$; no other short modulus is needed.
The remaining tuples lie in $\mathcal R_{\rm bad}(\eta_1)$ up to
$O(1/\cL)$. Use nonnegative smooth cells of bounded overlap supported in
$\mathcal R_{\rm bad}(2\eta_1)$.  Write
\begin{equation}\label{C:eq:cell-scales}
 \begin{gathered}
 c=N^\gamma,\quad a\asymp A=N^\alpha,\quad d\asymp D,\quad AD\asymp N/c,\\
 e\asymp E=N^{\beta-\gamma},\qquad f\asymp F=N^{1+\alpha-\beta}.
 \end{gathered}
\end{equation}
Cells have side $O(1/\cL)$ in the exponent variables.  The spectral
weights may instead be chosen in $(f',a/f')$, $f'\in\{e,f\}$;
these coordinates give the same bounded-overlap covering.  Positive
majorants allow us to drop the original size restrictions inside a cell.
Proposition~\ref{M:3.3} gives, throughout the enlarged bad region,
\begin{equation}\label{C:eq:large-divisors}
 e,f\gg N^{1/2-3\eta_1}.
\end{equation}
For example $\beta\geq(1-2\eta_1)/2+O(1/\cL)$, while
$1+\alpha-\beta\geq
\max(1-\alpha-\gamma-2\eta_1,\alpha-2\eta_1)+O(1/\cL)$.

\smallskip\noindent
\emph{2. Bands and layers.}
Let $c\asymp2^j$, $0\leq j\ll\eta\cL$, and put
\[
 \delta=|2\alpha-1+\gamma|,\qquad k=\lfloor\delta\cL\rfloor.
\]
A layer means fixed $(j,k)$ and all its $\beta$-cells; there are only
$O(1)$ choices of the $a$-block for each $c$-block and $k$.
Fixed shifts in $j,k$ and fixed changes of logarithm base are absorbed
in constants.  Remove bands of width $C_0/\cL$ about
$\beta=\alpha+\gamma$ and $\beta=1$, and also the cells with $D\leq64$.
In each such cell one of $e,f$ is $O_{C_0}(A)$ (for the last band use
$ef=4a^2d+1$).  Lemma~\ref{L:8.4}(b), with a fixed enlarged constant
in $F\leq4A$, bounds its mass by $O(2^jAD)=O(N)$.
There are $O(\cL)$ band cells per $c$-block, so all bands cost
$O(N\cL^2)$ without a sieve.  Layers $k\leq C_1\log\cL$ are postponed.

For reference, Brun--Titchmarsh on $4ad$ bounds a whole layer by
\begin{equation}\label{C:eq:BT-layer}
 O(N\cL/j),\qquad j\geq1,
\end{equation}
whenever $k\leq\cL/3$.  For fixed $(a,d,f)$ the $c$-block maps into
an interval of length $4ad\,2^j$ in one residue class modulo $4ad$.
Brun--Titchmarsh in the interval form
$\pi(x+y;m,b)-\pi(x;m,b)\le2y/(\varphi(m)\log(y/m))$ \cite{MV}
(R114 repair: not \cite[(A.10)]{ET}, which concerns initial segments),
with $m=4ad$ and $y=4ad\,2^j$, gives
$8ad\,2^j/(\varphi(4ad)j\log2)$ for a reduced class.  In a
nonreduced class there is at most one prime, which is absorbed by the
same upper bound.  Now $k\leq\cL/3$ implies $A,D\geq N^{1/4}$ for
our fixed small $\eta$, and Lemma~\ref{L:8.4}(a2), summed over the
divisors $f$, proves \eqref{C:eq:BT-layer}.  For $j=0$ use its trivial
analogue $O(N\cL)$.

\smallskip\noindent
\emph{3. The cell estimates away from the bands.}
All errors in the next paragraphs are \emph{summed absolute errors}
normalized by $AD$ (by $N$ after the $c$-sum), using Cauchy--Schwarz
and Lemma~\ref{L:6.3} for Theorem~\ref{L:6.2}. The bad region gives
$A\ge N^{1/3-2\eta_1}$, hence $A\ge\cL^2$ for Lemma~\ref{L:8.4}(a1).
All spectral boxes satisfy $Y/\lambda\asymp(A\sqrt d)^{-1}
\ll N^{-1/3+2\eta_1}$ and $L_*\ll\cL$, so Proposition~\ref{L:5.1}'s
exact height condition holds for large $N$. We write $F'$ for the scale of the selected divisor.

\emph{(i) $D\geq A$.}
Use the fixed-$a$ sequences of Proposition~\ref{L:7.1}.
By \eqref{C:eq:Ka-summed} and \eqref{C:eq:large-divisors} their
normalized error, summed with sieve weights through $q\leq Q$, is
\[
 \ll\cL^C\bigl(Q^2N^{-\delta}+QN^{-c_0}\bigr),
 \qquad c_0=\tfrac12-4\eta_1>0.
\]
There is no restriction on $\beta$.  The bad region gives
$\alpha\geq1/3-2\eta_1$, hence $\delta\leq1/3+4\eta_1$ here.
Thus $Q\leq N^{\delta/4}$ makes the second term an absolute power
saving, for sufficiently small $\eta_1$.

\emph{(ii) $D<A$ and $\beta<\alpha+\gamma$.}
Here $e\ll A$.  If $k\geq2j$ (allowing fixed shifts of the thresholds),
use the $e$-cusp variant of Theorem~\ref{L:6.2}, whose density is
unchanged by the involution proved in \S\ref{sec:perd}. Thus $F'=E$ and
$\lambda\asymp A/(qE)$, including $\lambda>1$.
Its summed error is
\begin{equation}\label{C:eq:e-cusp}
 \ll\cL^C Q^2
 \left[N^{-\delta/2}(1+A/E)^{1/2}
       +\frac{E^{1/2+\varepsilon}}A\right].
\end{equation}
Since $b\geq a/2$, we have $\beta\geq\alpha-O(1/\cL)$ and
$(A/E)^{1/2}\ll N^{\gamma/2}\ll N^{\delta/4}$ in this range.
The other term is $\ll A^{-1/2+2\varepsilon}$ and is
$O(N^{-\delta/4})$ on the bad region, for a fixed sufficiently small
$\varepsilon$.  This gives $\cL^CQ^2N^{-\delta/4}$.
Cells with $a/2\leq b<a$ are included: they satisfy $e\leq3a/c$.
If $k<2j$, use \eqref{C:eq:BT-layer} instead; its hypotheses hold
because $j\ll\eta\cL$.

\emph{(iii) $D<A$ and $\alpha+\gamma<\beta<1$.}
Off the bands, both $e,f\geq8A$. Choose the smaller divisor,
with scale $F'=\min(E,F)\ll A\sqrt D$, as cusp parameter.  Then
$\lambda\ll1/q$, and Theorem~\ref{L:6.2} gives
\[
 \ll\cL^C Q^2
 \left[N^{-\delta/2}+N^\varepsilon D^{1/4}A^{-1/2}\right].
\]
The second exponent is $(1-\gamma-3\alpha)/4+\varepsilon$;
its difference from $-\delta/4$ is $-\alpha/4+\varepsilon$.
It is therefore bounded by $N^{-\delta/4}$ with an absolute margin.
We use the weaker common error $\cL^CQ^2N^{-\delta/4}$.

\emph{(iv) $D<A$ and $1<\beta\leq1+2\eta_1$.}
This strip cannot be discarded: here $f<A$, even though $e$ is large.
Put $k'=\lfloor(\beta-1)\cL\rfloor$, which exceeds a fixed constant
off the band.  If $k'\leq k/2$, use the $f$-cusp, not the $e$-cusp.
With $F'=F$ and $\lambda\asymp A/(qF)$, Theorem~\ref{L:6.2} gives
\[
 \ll\cL^CQ^2
 \left[N^{-\delta/2}(1+A/F)^{1/2}+F^{1/2+\varepsilon}/A\right]
 \ll\cL^CQ^2N^{-\delta/4},
\]
because $(A/F)^{1/2}\ll N^{(\beta-1)/2}\ll N^{\delta/4}$.
The cusp term has an absolute margin.  Crucially, the first term
carries no $N^\varepsilon$ loss at the degenerating boundary.
If $k'>k/2$, use the $cdf$ fibres of Lemma~\ref{L:8.2}:
$4cdf\asymp N^{2-\beta}$ and $\log(N/(4cdf))\gg k'$.
Brun--Titchmarsh costs $O(N/k')$ per cell.  Indeed its reciprocal-totient
sum is bounded by
\[
 \sum_{c\asymp2^j}\frac1{\varphi(c)}
 \sum_{d\asymp D}\frac1{\varphi(d)}
 \sum_{f\asymp F}\frac{\rho_d(f)}{\varphi(f)}\ll1
\]
by Lemma~\ref{L:8.3}(c). Comparing with this weighted sum bounds the
nonreduced-class errors by $\sum_{c,d,f}\rho_d(f)\ll 2^jDF
\asymp N\exp(-k')\ll N/k'$. On the sieve side the weighted cell mass is $O(N)$
by Lemma~\ref{L:8.4}(b), since $F\leq A$.  Thus in either treatment
the cell costs $O(N/\max(k,k'))$ once the sieve saving below is established.

These cases cover every remaining cell: on $D\geq A$ use (i), and on
$D<A$ the three intervals separated by $\alpha+\gamma$ and $1$ give
(ii)--(iv).  Their boundaries and $A\asymp D$ were already set aside.

\smallskip\noindent
\emph{4. Obtaining the prime saving from the cell errors.}
For a sequence $\sigma$ with fixed $s\in\{a,d\}$ and fixed $c$, choose
\[
 z=N^{\kappa\delta/8},\qquad Q=z^2,
 \quad \kappa=1\text{ in (i)},\quad \kappa=1/4\text{ in (ii)--(iv)}.
\]
Sieve only by primes $\ell>\ell_0\geq3$ with $\ell\nmid2cs$.
Thus $(q,2ca)=1$ in Proposition~\ref{L:7.1}, and $(q,2cd)=1$ in
Theorem~\ref{L:6.2}; in particular $\kappa_c(q)=1$. Both densities satisfy
\[
 \frac1\ell-\frac2{\ell^2}\leq g_\sigma(\ell)\leq\frac1\ell<1.
\]
Selberg's upper-bound sieve \cite[Theorem~4.1 and Lemma~4.1]{HR}
for nonnegative weighted sequences therefore gives
\begin{equation}\label{C:eq:sieve}
 \begin{aligned}
 \#\{n\in\sigma:n\text{ prime},\ n>N/2\}
 &\leq\frac{X_\sigma}{G_\sigma(z)}+
 \sum_{\substack{q\leq Q\\q\mid P_\sigma(z)}}
 3^{\omega(q)}|r_\sigma(q)|,\\
 G_\sigma(z)&\gg\frac{\varphi(2cs)}{2cs}\log z.
 \end{aligned}
\end{equation}
Here counts are weighted if the cell is smooth; $z<N/2$.
The lower bound on the local density, not merely its upper bound,
is used in the estimate for $G_\sigma$.

The model mass $X_\sigma$ need not be compared from below to the actual
mass.  At $q=1$, Theorem~\ref{L:6.2} or Proposition~\ref{L:7.1} gives
\[
 X_\sigma\leq\#\sigma+|r_\sigma(1)|.
\]
Also $2cs/\varphi(2cs)\leq2g(c)g(s)$.
Lemma~\ref{L:8.4}(a1) and $\sum_{c\asymp2^j}g(c)\ll2^j$ bound
the weighted mass of a layer in (i) or (iii) by $O(N\cL)$.
Lemma~\ref{L:8.4}(b) bounds each cell in (ii) or (iv) by $O(N)$;
in particular their layer masses are $O(N\cL)$.
The $q=1$ errors introduced in this comparison have only the extra
factor $\max g(c)g(s)\ll(\log\cL)^2$.
Since $Q^2=N^{\kappa\delta/2}$, the summed errors are
$\ll\cL^CN^{-\kappa\delta/2}$ times the relevant mass scale;
the additional term in (i) has an absolute power saving.
All these errors are $O(1/k)$ times that scale once
$k\geq C_1\log\cL$, with $C_1$ fixed sufficiently large.
The main term in \eqref{C:eq:sieve} has the same saving, since
$\log z\asymp k$.  This proves all the sieve savings asserted in step~3.

\smallskip\noindent
\emph{5. Summing all cells.}
For $1\leq j\ll\eta\cL$, the layers $k\leq C_1\log\cL$ cost,
by \eqref{C:eq:BT-layer},
\[
 \ll N\cL\log\cL\sum_{j\ll\eta\cL}\frac1j
 \ll N\cL(\log\cL)^2\ll N\cL^2.
\]
For $j=0$ these layers cost $O(N\cL\log\cL)$ trivially.
In (i), (iii), and (ii) with $k\geq2j$, a layer costs
\[
 \ll N\cL\min(1/j,C/k)\qquad(k\leq\cL/3).
\]
One takes the better bound for the whole portion of the layer in
question, not an unweighted average of the two methods.
Lemma~\ref{L:1.1} sums these costs to $O(N\cL^2)$.
For $k>\cL/3$ the sieve alone suffices:
$\sum_j\sum_{k>\cL/3}N\cL/k\ll N\cL^2$.
Case (ii) with $k<2j$ costs
$\sum_j\sum_{k<2j}O(N\cL/j)=O(N\cL^2)$.
In (iv) the per-cell bound gives, for each $c$-block,
\[
 \sum_{k,k'\ll\cL}\frac{CN}{\max(k,k')}
 \ll N\sum_{h\ll\cL}\frac{2h-1}{h}\ll N\cL.
\]
Summing over $j$ gives $O(N\cL^2)$ again.  For $j=0$ outside the
innermost layers the $1/k$ bound costs at most $O(N\cL\log\cL)$,
and the preceding two-dimensional sum still applies to (iv).
Together with the bands and step~1, this proves the bound on
$(N/2,N]$.  Summing dyadically in $N$ proves the theorem.
\end{proof}

\section{What is unconditional}\label{sec:uncond}
The separation argument, Sobolev duality, the local-density computation,
the per-$a$ count, and the arithmetic mass estimates are unconditional.
The use of (SEL) is confined to the exceptional spectrum in
Proposition~\ref{L:5.1}.  It is nevertheless essential to the present
proof of the prime-counting theorem.

\subsection{The exceptional-spectrum obstruction}
An eigenvalue $\lambda_j=1/4-\sigma_j^2$, $0<\sigma_j<1/2$, introduces
in the unfolded coefficient the factor $(nY)^{-\sigma_j}$, since
$K_{\sigma_j}(2\pi ny)$ grows like a constant times $(ny)^{-\sigma_j}$
when $ny$ is small.  The unconditional Kim--Sarnak bound
\cite[Appendix~2]{KS} is $\sigma_j\leq7/64$, including the
nebentypus forms at issue here.  A direct use of this bound permits the
loss
\begin{equation}\label{C:eq:exceptional-loss}
 Y^{-7/64}\asymp(qF'\sqrt d)^{7/64}
\end{equation}
in the spectral error.  On the $D<A$ side this must be compared with
$(D/A)^{1/2}=N^{-\delta/2}$, $\delta=2\alpha-1+\gamma>0$.
In the worst cells $qF'\sqrt d$ has size about $N$, so the exponent
comparison leaves a strip
\[
 0<2\alpha-1+\gamma\lesssim7/32.
\]
Even the smallest cusp parameters at the boundary, with
$F'\gg N^{1/2-3\eta_1}$ and $D$ near $N^{1/2}$, give a loss of
size $N^{21/256-O(\eta_1)}$, corresponding to the comparison
$\delta\approx21/128-O(\eta_1)$.  These numbers describe the
obstruction to the estimates, not a sharp geometric boundary or a
lower bound for an actual spectral contribution.  In particular a
fixed positive width, rather than a single line, remains untreated.

\begin{assessment}[partial exceptional-spectrum estimates]\label{C:ass:exceptional}
Deshouillers--Iwaniec's exceptional large sieve
\cite[Theorem~5]{DI} improves on a pointwise insertion of
\eqref{C:eq:exceptional-loss}, but does not supply the needed uniform
saving at $A\asymp D$.  On a Fourier block $n\asymp N_0$ its parameter
must be
\[
 X=\frac1{N_0Y},
\]
not $1/Y$, because the squared exceptional weight is
$(nY)^{-2\sigma_j}$.  In its usual notation the bound has the shape
\[
 \sum_{\rm exc}X^{2\sigma_j}
 \left|\sum_{n\asymp N_0}a_n\rho_j(n)\right|^2
 \ll (1+\sqrt{\mu N_0X})(1+\sqrt{\mu N_0^{1+\varepsilon}})
 \sum_n|a_n|^2,
\]
with $\mu\asymp1/M$, $M=4dq^2$.  The first factor is of size
$1+(F'/(q\sqrt d))^{1/2}$, which is not uniformly harmless near the
boundary.  The level average in \cite[Theorem~6]{DI} offers further
partial control.  The source calculation reports an extra factor
$(F')^{3/4}D^{1/8}/A$ there; its detailed bookkeeping was not re-checked
in the internal review, and we do not use it as a proved obstruction
theorem.  The same caveat applies to the source's exploratory
exceptional-density calculations.  None enters the conditional proof.
\end{assessment}

A positive-width region bounded only by the original $1/j$ saving
still contributes at the level of the available upper bounds
$N\cL^2\sum_{j\ll\cL}1/j$.  Unlike the linearly degenerating
saving in Lemma~\ref{L:1.1}, this does not remove $\log\log N$.
Thus \emph{this argument gives no unconditional improvement} of
\cite[Theorem~1.1]{ET}; this is not a claim that an unconditional
improvement by other methods is impossible.  Sufficient replacements
for (SEL) would include a large sieve for the exceptional spectrum
with precisely these Fourier weights and $X\asymp1/(N_0Y)$, with
errors that still allow the layer summation, or cancellation beyond
termwise Weil in $\sum_{a\asymp A}S(h,k;4a^2)$.

\subsection{Lifts versus equidistribution of a finite packet}
The lift-counting target (H**) of \cite[\S3.5]{MN3} is stronger than
what the proof above uses.  For the packets of natural class-number
size $d^{1/2+o(1)}$ on a surface of volume $d^{1+o(1)}$, a box of
hyperbolic area $V$ has expected packet count $Vd^{-1/2+o(1)}$.
Consequently a uniform relative equidistribution statement in
\emph{all} boxes cannot extend below area $d^{1/2-o(1)}$: integrality
and boxes containing a packet point already prevent it.  Here
``scale'' is area, not hyperbolic radius.  This elementary granularity
obstruction does \emph{not} refute (H**) as formulated for lifts,
averaged in $d$, with congruence conditions.  Nor does it obstruct the
second-moment method here: the lift-counting Poincar\'e function has
box area $\asymp A\sqrt d$, and Proposition~\ref{L:4.3} uses only its
variance against a separated packet.

For comparison, Liu--Masri--Young
\cite[Theorem~1.4, and the more general Theorem~7.1]{LMY} prove a
joint level--discriminant equidistribution result in the range
$q\leq |D|^{1/20-\varepsilon}$.  Their $q$ is a prime level, $-D$ is
an odd fundamental discriminant, and $q$ splits in
$\Q(\sqrt{-D})$.  Their count is in translates of a level-one
fundamental domain; it is not the present thin-box count with sieve
twists.  Here level and discriminant are tied, already $d$ and $-4d$
before inserting $q$, and the coprimality/splitting hypotheses do not
match.  Thus that theorem does not supply (H**) or the needed per-$d$
estimate.  We make no exhaustive literature claim that no other
joint-aspect theorem exists.

\section{Open problems}\label{sec:open}
\begin{problem}[unconditional removal]
Prove $\sum_{p\leq N}f_I(p)\ll N\log^2N$ without (SEL), or obtain an
unconditional improvement of its present $\log\log N$ loss.  A proof
must address a positive-width strip adjacent to $A=D$, not merely
the line $A=D$ itself.
\end{problem}

\begin{problem}[weighted exceptional density]
Prove a density or large-sieve theorem for exceptional eigenvalues
on $\Gamma_0(4dq^2)$, with even nebentypus modulo $2q$, incorporating
the coefficients of the box Poincar\'e series and the weights
$(N_0Y)^{-2\sigma_j}$.  The relevant target is an averaged error that
retains a saving degenerating at most linearly in
$|2\alpha-1+\gamma|$, after summation in $d$ and the sieve modulus;
a count of exceptional eigenvalues alone need not suffice.
\end{problem}

\begin{problem}[square-modulus Kloosterman sums]
Obtain useful cancellation in $\sum_{a\asymp A}S(h,k;a^2)$, or its
$4a^2$ variant, uniformly in the Fourier ranges produced by
\eqref{C:eq:poisson}.  For the application the averaging must handle
the accompanying smooth weights, gcd factors, and the sieve
congruence.  Termwise square-root cancellation is exactly the
estimate that stops at $D\asymp A$.
\end{problem}

\begin{problem}[joint level--discriminant counting]
Establish a version of the lift-counting target (H**) at level
comparable to the discriminant, with the parity restriction and
congruence twists retained.  A second-moment theorem for the
appropriate Poincar\'e functions could suffice and would avoid a
false assertion of uniform fine-scale equidistribution of a finite
packet.
\end{problem}

\begin{problem}[general numerators]
Extend the conditional Type~I estimate uniformly from $4/n$ to
$m/n$.  Primes dividing $m$ must be omitted from the sieve, and the
mass gain $\varphi(m)/m$ must be tracked against the sieve-product
inflation $m/\varphi(m)$ and all remainder terms
\cite[remarks after Theorem~3.8]{MN3}.  The earlier internal transition
gap $(m\log m/\varphi(m))^{1/3}$ is historical: the refined lower-side
bound in \cite[\S2, Theorem~L$'$]{MN3} reduces it to
$(\log\log N)^{1/3}$.  An $m$-uniform removal of the Type~I logarithmic
loss would close that remaining gap at the scale
$\log N\asymp m^{1/3}$.  Such an extension is not proved here.
\end{problem}

\begin{remark}[numerical status: \status{evidence}]
The script \texttt{scripts/ttl\_perd.py} supplies only a weak sanity
check for per-$d$ counts.  Its weights are normalized to mass about
$20$, and its local-density comparison is not the orbital main term
of Theorem~\ref{L:6.2}.  Small observed errors are therefore neither
evidence for a uniform spectral error bound at the required scales
nor a test of (SEL).  No numerical observation is used in any proof.
\end{remark}
\begin{thebibliography}{99}
\bibitem[ET]{ET} C. Elsholtz and T. Tao, Counting the number of solutions to the Erd\H{o}s--Straus equation on unit fractions, \emph{J. Aust. Math. Soc.} \textbf{94} (2013), 50--105; arXiv:1107.1010v6. Page and equation references in this note refer to that arXiv version.
\bibitem[MN3]{MN3} \emph{The lower side of the $m/n$ transition: the $m/\varphi(m)$ loss and the Type I $\log\log$}, internal note, \texttt{EXCEPTIONAL\_MN3.md}, with internal review \texttt{reviews/exceptional-mn3-review.md}.
\bibitem[DI]{DI} J.-M. Deshouillers and H. Iwaniec, Kloosterman sums and Fourier coefficients of cusp forms, \emph{Invent. Math.} \textbf{70} (1982), 219--288.
\bibitem[Dr]{Dr} S. Drappeau, Sums of Kloosterman sums in arithmetic progressions, and the error term in the dispersion method, \emph{Proc. Lond. Math. Soc.} \textbf{114} (2017), 684--732; arXiv:1504.05549. Proposition numbers here use the arXiv version.
\bibitem[IK]{IK} H. Iwaniec and E. Kowalski, \emph{Analytic number theory}, American Mathematical Society Colloquium Publications \textbf{53}, American Mathematical Society, 2004.
\bibitem[MV]{MV} H. L. Montgomery and R. C. Vaughan, The large sieve, \emph{Mathematika} \textbf{20} (1973), 119--134.
\bibitem[Jia]{Jia} C. Jia, The estimate for mean values on prime numbers relative to $4/p=1/n_1+1/n_2+1/n_3$, \emph{Sci. China Math.} \textbf{55} (2012), 465--474.
\bibitem[LMY]{LMY} S.-C. Liu, R. Masri and M. P. Young, Subconvexity and equidistribution of Heegner points in the level aspect, \emph{Compos. Math.} \textbf{149} (2013), 1150--1174; arXiv:1206.3208.
\bibitem[Duke]{Duke} W. Duke, Hyperbolic distribution problems and half-integral weight Maass forms, \emph{Invent. Math.} \textbf{92} (1988), 73--90.

\bibitem[Iw]{Iw} H. Iwaniec,
\emph{Spectral Methods of Automorphic Forms}, second edition,
Graduate Studies in Mathematics \textbf{53}, American Mathematical Society, 2002.

\bibitem[HR]{HR} H. Halberstam and H.-E. Richert,
\emph{Sieve Methods}, Academic Press, London, 1974.

\bibitem[KS]{KS} H. H. Kim and P. Sarnak,
Refined estimates towards the Ramanujan and Selberg conjectures,
Appendix~2 to H. H. Kim, Functoriality for the exterior square of
$\mathrm{GL}_4$ and the symmetric fourth of $\mathrm{GL}_2$,
\emph{J. Amer. Math. Soc.} \textbf{16} (2003), 139--183.

\bibitem[Weil]{Weil} A. Weil,
On some exponential sums,
\emph{Proc. Natl. Acad. Sci. USA} \textbf{34} (1948), 204--207.

\end{thebibliography}
\end{document}
